%% toc: P2 21 21.4,21.5,21.6,21.7
\section{Proofs for the decision procedures of the decidable fragments}
\label{sec:annex-p02-w02-p02-035-decision-procedures}

Throughout, structures have nonempty domains, a formula is evaluated under an
assignment of elements to its free variables, and a quantifier-free formula is
true under an assignment exactly when the propositional valuation of its
atoms induced by the assignment makes it true.  Two counting facts are used
repeatedly.  First, a code of length \(s\) over a fixed alphabet with \(b\geq2\)
letters contains \(O(s/\log s)\) distinct symbol names: at most \(b^{l+1}\)
names have length at most \(l\), so among \(r\) distinct names at least \(r/2\)
have length at least \(\log_b(r/2)-1\), whence
\(s\geq(r/2)(\log_b(r/2)-1)\).  Second, a nondeterministic machine that guesses
a string of length \(L\) and then checks it in time \(T\) runs in time
\(O(L+T)\).

\begin{annexitem}{annex:p02-w02-p02-035-domino-hard}{thm:p02-w02-p02-035-domino-hard}{The bounded torus domino problem is NEXPTIME-hard}
\begin{proof}
Let \(L\in\mathsf{NEXPTIME}\), say \(L\in\mathsf{NTIME}(2^{n^k})\) via a
multitape machine.  By premise (i) of
\cref{def:p02-w02-p02-035-domino-import}, \(L\) is accepted by a one-tape
machine within time \(C2^{2n^k}\) for some constant \(C\).  Choose \(j\) with
\(\log_2C+2n^k\leq(n+2)^j\) for all \(n\geq0\), and put
\(T(n)\defeq2^{(n+2)^j}\).  Then \(L\) is accepted by a one-tape machine in
time \(T(n)\), hence in space \(T(n)+1\), and \(T(n)\geq2n\).  By premise
(ii) there is a simple machine \(M\) accepting \(L\) within these bounds.

\emph{Padding.}  Let \(\natural\) be a symbol not in the input alphabet
\(\Sigma\) of \(M\).  For a constant \(d\geq1\), fixed below, put
\(N(n)\defeq(n+2)^j+d\) and
\[
 L^\natural\defeq\{x\natural^{N(\lvert x\rvert)-\lvert x\rvert}:x\in L\}.
\]
A one-tape machine \(M^\natural\) accepts \(L^\natural\) as follows: it checks
deterministically that its input \(u\) has the form
\(x\natural^{N(\lvert x\rvert)-\lvert x\rvert}\), erases the padding, returns
the head to cell \(0\) and runs \(M\) on \(x\).  There is a nondecreasing
polynomial \(q\) with \(q(l)\geq2l\) such that the first three phases take
time and space at most \(q(\lvert u\rvert)\).  So \(M^\natural\) accepts each
\(u\in L^\natural\), \(\lvert u\rvert=N\), \(\lvert x\rvert=n\), within time
\(t_0\leq q(N)+T(n)\) and space \(s_0\leq t_0+1\); and it accepts no
\(u\notin L^\natural\).  By premise (ii) we may take \(M^\natural\) simple
with the same bounds.

Choose \(d\geq1\) such that \(2^{d-1}\geq q(2d)+3\) and \(2^{l}\geq q(2l)+3\)
for every \(l\geq d\); such \(d\) exists because \(2^l\) eventually exceeds
every polynomial in \(l\).  Write \(t=(n+2)^j\geq1\), so \(N=t+d\).  If
\(t\leq d\), then \((2^d-1)2^t\geq2^{d-1}\geq q(2d)+3\geq q(t+d)+3\).  If
\(t\geq d\), then \((2^d-1)2^t\geq2^t\geq q(2t)+3\geq q(t+d)+3\).  In both
cases
\[
 2^N=2^t+(2^d-1)2^t\geq T(n)+q(N)+3\geq t_0+3\geq s_0+2 .
\]

\emph{Tiling.}  Premise (iii) applied to \(M^\natural\) gives a domino system
\(\mathcal D\) and a linear-time map \(u\mapsto w\) with
\(\lvert w\rvert=\lvert u\rvert\).  Let
\(g(x)\) be the word \(w\) obtained from
\(u=x\natural^{N(\lvert x\rvert)-\lvert x\rvert}\); \(g\) is computable in
polynomial time and \(\lvert g(x)\rvert=N\).  If \(x\in L\), then
\(M^\natural\) accepts \(u\) within time \(t_0\) and space \(s_0\), and
\(2^N\geq t_0+2\), \(2^N\geq s_0+2\); so \(\mathcal D\) tiles
\(U(2^N,2^N)=Z(2^N)\) with initial condition \(w\), that is,
\(w\in\mathrm{DOMINO}(\mathcal D,2^n)\).  If \(x\notin L\), then \(u\notin
L^\natural\), \(M^\natural\) does not accept \(u\), and \(\mathcal D\) tiles no
\(U(s,t)\) with \(s,t\geq2\) with initial condition \(w\); in particular
\(w\notin\mathrm{DOMINO}(\mathcal D,2^n)\), since \(2^N\geq2\).
\end{proof}
\end{annexitem}

\begin{annexitem}{annex:p02-w02-p02-035-monadic-small-model}{thm:p02-w02-p02-035-monadic-small-model}{Colour collapse and the \(2^m\) bound}
\begin{proof}
The map \(c:A\to A_{\mathrm{col}}\) is surjective by the definition of
\(A_{\mathrm{col}}\), and \(\mathfrak A\entails P_ia\) iff \(c(a)_i=1\) iff
\(\mathfrak A_{\mathrm{col}}\entails P_i\,c(a)\).  For an assignment \(\sigma\)
into \(A\) write \(c\circ\sigma\) for the assignment into
\(A_{\mathrm{col}}\).  We prove by induction on \(\varphi\) that
\(\mathfrak A\entails\varphi[\sigma]\) iff
\(\mathfrak A_{\mathrm{col}}\entails\varphi[c\circ\sigma]\) for every
\(\sigma\).

Atomic formulas are \(P_ix\), since there is no equality and no symbol of
other arity; the claim is the displayed equivalence with \(a=\sigma(x)\).
Negation and conjunction preserve the claim because both sides use the same
truth tables.  Let \(\varphi=\exists y\chi\).  If
\(\mathfrak A\entails\chi[\sigma[y\mapsto a]]\), then by induction
\(\mathfrak A_{\mathrm{col}}\entails\chi[(c\circ\sigma)[y\mapsto c(a)]]\),
because \(c\circ(\sigma[y\mapsto a])=(c\circ\sigma)[y\mapsto c(a)]\).
Conversely, if \(\mathfrak A_{\mathrm{col}}\entails
\chi[(c\circ\sigma)[y\mapsto e]]\) for some colour \(e\), surjectivity gives
\(a\) with \(c(a)=e\), and the same identity with induction gives
\(\mathfrak A\entails\chi[\sigma[y\mapsto a]]\).  The universal quantifier is
\(\neg\exists\neg\).

For a sentence \(\psi\) the claim gives \(\mathfrak A\entails\psi\) iff
\(\mathfrak A_{\mathrm{col}}\entails\psi\).  The universe of
\(\mathfrak A_{\mathrm{col}}\) is a nonempty subset of \(\{0,1\}^m\), so it has
at most \(2^m\) elements.
\end{proof}
\end{annexitem}

\begin{annexitem}{annex:p02-w02-p02-035-monadic-procedure}{thm:p02-w02-p02-035-monadic-procedure}{Correctness and an exponential-time bound}
\begin{proof}
\emph{Correctness.}  If Step 3 succeeds for some \(S\), then
\(\mathfrak B_S\) is a model and \(\psi\) is satisfiable.  If
\(\mathfrak A\entails\psi\), let \(S\) be the set of colours realized in
\(\mathfrak A\).  Then \(\mathfrak B_S=\mathfrak A_{\mathrm{col}}\), and
\cref{thm:p02-w02-p02-035-monadic-small-model} gives
\(\mathfrak B_S\entails\psi\).

\emph{Time.}  Each predicate letter occurs in \(\psi\), so \(m\leq n\).  The
guessed set \(S\) is written as at most \(2^m\) words of length \(m\), at most
\(n2^n\) symbols.  Evaluate by the recursive truth definition, with an
assignment stored as a list of colours.  Let \(N=\lvert S\rvert\leq2^n\).  For
a subformula \(\chi\) of length \(l\) and quantifier depth \(e\), the number
of recursive calls made in evaluating \(\chi\) under one assignment is at most
\(lN^{e}\): this holds for atoms; for a Boolean combination it follows by
adding the bounds for the parts; and for \(\exists y\chi'\) there are \(N\)
evaluations of \(\chi'\), giving at most \(1+N(l-1)N^{e-1}\leq lN^e\).  Since
\(l,e\leq n\), there are at most \(n2^{n^2}\) calls.  Each call reads an
assignment of at most \(n\) colours and looks up colours in \(S\), which takes
time polynomial in \(n2^n\).  The total time is \(2^{O(n^2)}\), which lies in
\(\mathsf{NTIME}(2^{n^3})\subseteq\mathsf{NEXPTIME}\) for large \(n\).
Validity of \(\psi\) is non-satisfiability of \(\neg\psi\), so it is in
\(\mathsf{coNEXPTIME}\).
\end{proof}
\end{annexitem}

\begin{annexitem}{annex:p02-w02-p02-035-monadic-hard}{thm:p02-w02-p02-035-monadic-hard}{Monadic satisfiability is NEXPTIME-complete}
\begin{proof}
Write \(\sigma\) for the successor on \(Z(2^n)\) in row-major order:
\(\sigma(x,y)=(x+1,y)\) if \(x<2^n-1\) and \(\sigma(2^n-1,y)=(0,y+1)\), with
\(y+1\) taken modulo \(2^n\).  It adds \(1\) modulo \(2^{2n}\) to the number
\(x+2^ny\), so it is a cyclic permutation of all \(2^{2n}\) points.  Binary
increment modulo \(2^n\) flips bit \(i\) exactly when all lower bits are \(1\):
if \(x'=x+1\bmod2^n\), then \(x'_i=x_i\oplus\bigwedge_{j<i}x_j\).

\emph{From a tiling to a model.}  Let \(\tau\) witness
\(w\in\mathrm{DOMINO}(\mathcal D,2^n)\).  Interpret the predicates on
\(Z(2^n)\) as intended, with \(N_i=\{(i,0)\}\) for \(i\leq n\) (these points
exist and are distinct since \(n<2^n\)), and \(P_d=\tau^{-1}(d)\).  For each
\(z\) take \(v=\sigma(z)\).  The \(X^{*}\) and \(Y^{*}\) clauses of \(\alpha\)
hold by the definition of \(X_i^{*},Y_i^{*}\).  The \(X\)-clause is the
increment rule for \(x\).  The \(Y\)-clause says that \(y\) is incremented
exactly when \(X_{n-1}z\land X_{n-1}^{*}z\), that is, when all bits of \(x\)
are \(1\), which is the definition of \(\sigma\).  \(N_0z\) holds iff all bits
are \(0\).  For \(i<n\), \(z=(i,0)\) iff \(\sigma(z)=(i+1,0)\), because
\(\sigma\) is injective and \(\sigma(i,0)=(i+1,0)\) as \(i<2^n-1\).  So
\(\forall z\exists v\alpha\) holds.  For \(\beta\): exactly one \(P_dz\) holds
since \(\tau\) is a function.  By the increment rule, \(H(z,z')\) holds iff
\(z'=(x+1\bmod2^n,y)\) and \(V(z,z')\) iff \(z'=(x,y+1\bmod2^n)\); the tiling
conditions give the two implications.  \(N_iz\) means \(z=(i,0)\), where
\(\tau(i,0)=w_i\).

\emph{From a model to a tiling.}  Let \(\mathfrak A\entails\psi_w\).  For
\(a\in A\) let \(f(a)=(x(a),y(a))\) be the point whose bits are the truth
values of \(X_0a,\ldots,X_{n-1}a\) and \(Y_0a,\ldots,Y_{n-1}a\).  By the first
three clauses of \(\alpha\), instantiated at \(z=a\), induction on \(i\) shows
that \(X_i^{*}a\) holds iff \(X_ja\) for all \(j<i\), and likewise for
\(Y^{*}\).  Then the next two clauses show that every witness \(v\) for
\(z=a\) satisfies \(f(v)=\sigma(f(a))\), and the definitions of \(H\) and
\(V\) show that \(H(a,b)\) holds iff \(f(b)\) is the right neighbour of
\(f(a)\), and \(V(a,b)\) iff \(f(b)\) is the upper neighbour.

Starting from any element and repeatedly taking witnesses, \(f\) runs
through the whole \(\sigma\)-orbit, which is \(Z(2^n)\); so \(f\) is
surjective.  Choose \(b_0\) with \(f(b_0)=(0,0)\); then \(N_0b_0\) holds.
Given \(b_i\) with \(i<n\), \(f(b_i)=(i,0)\) and \(N_ib_i\), let \(b_{i+1}\) be
a witness for \(z=b_i\); then \(f(b_{i+1})=(i+1,0)\), and the last clause of
\(\alpha\) gives \(N_{i+1}b_{i+1}\).  For each point \(z\) choose
\(a_z\in f^{-1}(z)\), taking \(a_{(i,0)}=b_i\) for \(i<n\), and let \(\tau(z)\)
be the unique \(d\) with \(P_da_z\).  For each \(z\), instantiating
\(\forall z\forall z'\beta\) at \((a_z,a_{z'})\) with \(z'\) the right
neighbour gives \(H(a_z,a_{z'})\) and so \((\tau(z),\tau(z'))\in H\); the
vertical case is the same.  Finally \(N_ib_i\) and the last clause of
\(\beta\) give \(\tau(i,0)=w_i\).  So \(\mathcal D\) tiles \(Z(2^n)\) with
initial condition \(w\).

\emph{Size.}  \(\psi_w\) has \(O(n)\) clauses of bounded shape in
\(\alpha\), \(H\) and \(V\), \(O(\lvert D\rvert^2)\) literals in the
exactly-one clause, and \(n\) initial clauses; each literal carries an index
of \(O(\log n)\) symbols.  So \(\lvert\psi_w\rvert=O(n\log n)\), and \(\psi_w\)
is produced from \(w\) by a polynomial-time loop.

\emph{Completeness.}  Membership is
\cref{thm:p02-w02-p02-035-monadic-procedure}.  For \(L\in\mathsf{NEXPTIME}\),
\cref{thm:p02-w02-p02-035-domino-hard} gives \(\mathcal D\) and \(g\); then
\(x\mapsto\psi_{g(x)}\) reduces \(L\) to monadic satisfiability in polynomial
time.  If \(L\in\mathsf{coNEXPTIME}\), the same construction for the
complement gives \(x\mapsto\psi_{g(x)}\) with \(x\in L\) iff \(\psi_{g(x)}\) is
unsatisfiable, iff \(\neg\psi_{g(x)}\) is valid.  Renaming \(z\) to \(x\) and
both \(v\) and \(z'\) to \(y\) puts \(\psi_w\) in the stated form; the renaming
does not change truth, because \(v\) and \(z'\) occur in different conjuncts.
\end{proof}
\end{annexitem}

\begin{annexitem}{annex:p02-w02-p02-035-scott}{thm:p02-w02-p02-035-scott}{Scott reduction}
\begin{proof}
\emph{Stage 1: arity at most two.}  Let \(\varphi\) have size \(s\) and
variables \(x,y\).  For each atomic subformula \(R(v_1,\ldots,v_k)\) with
\(k>2\), each \(v_l\in\{x,y\}\), introduce a symbol \(R_{(v_1\ldots v_k)}\),
named by the atom's own string with a marker so that its name is no longer
than a constant times the atom.  It is binary if both variables occur among
the \(v_l\), and unary otherwise; replace \(R(\bar v)\) by
\(R_{(\bar v)}(x,y)\), \(R_{(\bar v)}(x)\) or \(R_{(\bar v)}(y)\) accordingly.
Call the result \(\varphi^\dagger\).  Append: (a) for each such atom
\(R(\bar v)\) containing both variables, if the atom \(R(\bar w)\) obtained by
exchanging \(x\) and \(y\) also occurs in \(\varphi\), the sentence
\(\forall x\forall y(R_{(\bar v)}(x,y)\leftrightarrow R_{(\bar w)}(y,x))\); and
(b) for each such \(R\), the one-variable sentence stating, as a cycle of
implications, that all diagonal forms \(R_{(\bar v)}(x,x)\) and
\(R_{(\bar v)}(x)\) of its atoms are equivalent.  Let \(\varphi'\) be the
conjunction.  Every new conjunct has length linear in the atoms it mentions,
so \(\varphi'\) has size \(O(s)\), and it has at most \(s\) relation symbols.
This stage follows the source \cite[pp.~4--5]{GradelKolaitisVardi1997}.

If \(\mathfrak A\entails\varphi\), interpret \(R_{(\bar v)}\) at \((a,b)\) as
\(R^{\mathfrak A}(\bar v[x\mapsto a,y\mapsto b])\).  Atoms keep their truth
values, so \(\varphi^\dagger\) holds, and (a), (b) hold because both sides name
the same tuple.  Conversely let \(\mathfrak B\entails\varphi'\), and define
\(R^{\mathfrak A}\) on the universe of \(\mathfrak B\): if a tuple \(\bar c\)
equals \(\bar v[x\mapsto a,y\mapsto b]\) for an atom \(R(\bar v)\) of
\(\varphi\), let \(R^{\mathfrak A}(\bar c)\) be the value of the corresponding
new atom at \((a,b)\); otherwise let it be false.  This is well defined.  If
\(\bar c\) is constant, every representation is a diagonal form, and (b) makes
them agree.  If \(\bar c\) has two distinct entries \(a\neq b\), then any atom
representing it contains both variables, the positions holding \(a\) are
exactly the positions of one variable, and so the only representations are
\(\bar v\) at \((a,b)\) and its exchange \(\bar w\) at \((b,a)\), which (a)
makes agree.  Other symbols keep their interpretation.  Now every atom of
\(\varphi\) and its replacement have the same truth value under every
assignment, so induction on formulas gives \(\mathfrak A\entails\varphi\).
The universe is unchanged in both directions.

\emph{Stage 2: defining symbols.}  Write \(\varphi'\) with \(\land,\neg,\forall\)
only; this changes its size by a constant factor.  For each subformula
occurrence \(\chi\) with free variables \(\bar v\subseteq\{x,y\}\) introduce
\(Q_\chi\) of arity \(\lvert\bar v\rvert\in\{0,1,2\}\), with a binary index as
its name, and let \(\delta_\chi\) be
\(\forall\bar v(Q_\chi(\bar v)\leftrightarrow\chi^{\circ}(\bar v))\), where
\(\chi^{\circ}\) is \(\chi\) if \(\chi\) is atomic, \(Q_{\chi_1}\land
Q_{\chi_2}\) if \(\chi=\chi_1\land\chi_2\), \(\neg Q_{\chi_1}\) if
\(\chi=\neg\chi_1\), and \(\forall u\,Q_{\chi_1}\) if \(\chi=\forall u\chi_1\)
\cite[pp.~5--6]{GradelKolaitisVardi1997}.  Put
\(\varphi^{*}\defeq Q_{\varphi'}\land\bigwedge_\chi\delta_\chi\).

Each \(\delta_\chi\) is equivalent to a conjunction of sentences with prefix
\(\forall\forall\) or \(\forall\exists\); dummy quantifiers fill short
prefixes.  For atomic, conjunctive and negated \(\chi\) the matrix is
quantifier-free.  For \(\chi=\forall u\chi_1\) with \(t\) the other variable,
\(\delta_\chi\) is equivalent to
\[
 \forall t\forall u\bigl(Q_\chi(t)\to Q_{\chi_1}\bigr)\land
 \forall t\exists u\bigl(\neg Q_{\chi_1}\lor Q_\chi(t)\bigr),
\]
because \(u\) is not free in \(Q_\chi(t)\); when \(\chi\) is a sentence the
same holds with \(Q_\chi\) zero-ary.  Joining the \(\forall\forall\) conjuncts
gives the Scott form.

If \(\mathfrak A\entails\varphi'\), interpret each \(Q_\chi\) as the set of
tuples satisfying \(\chi\); every \(\delta_\chi\) holds and so does
\(Q_{\varphi'}\).  If \(\mathfrak B\entails\varphi^{*}\), induction on \(\chi\)
using \(\delta_\chi\) shows that \(Q_\chi\) and \(\chi\) agree on all tuples;
in particular \(Q_{\varphi'}\) gives \(\mathfrak B\entails\varphi'\).  Again
the universe is unchanged.  There are at most \(O(s)\) subformula occurrences,
each \(\delta_\chi\) has a bounded number of symbols besides at most three
names of length \(O(\log s)\) and, for atomic \(\chi\), a copy of the atom.  So
\(\varphi^{*}\) has size \(O(s\log s)\) and \(O(s)\) relation symbols, all of
arity at most \(2\).  Both stages are polynomial-time.

\emph{The \(x\neq y\) form.}  Let \(\mathfrak A\) have at least two elements
and \(\beta\) be quantifier-free.  If \(\mathfrak A\entails\beta(a,b)\) with
\(b\neq a\), then \(b\) witnesses \(x\neq y\land(\beta(x,x)\lor\beta(x,y))\)
at \(a\).  If \(b=a\), then \(\beta(a,a)\) holds and any \(b'\neq a\) is a
witness.  Conversely, a witness \(b\) for the new formula gives
\(\beta(a,a)\) or \(\beta(a,b)\), so \(\exists y\beta(a,y)\) holds.  The new
matrix entails \(x\neq y\).
\end{proof}
\end{annexitem}

\begin{annexitem}{annex:p02-w02-p02-035-scott-model}{thm:p02-w02-p02-035-scott-model}{Small models for the Scott class}
\begin{proof}
If \(m=0\), then \(\theta=\forall x\forall y\alpha\) is universal, and the
substructure of any model induced on one element is a model.  Let \(m\geq1\)
and \(\mathfrak A\entails\theta\).  Truth of a quantifier-free formula in two
elements depends only on their \(2\)-type, and in one element only on its
\(1\)-type; zero-ary symbols have the same value in every type of
\(\mathfrak A\), and all types below are copied from \(\mathfrak A\).  For a
\(2\)-type of \((a,b)\), the \(2\)-type of \((b,a)\) is obtained by exchanging
the variables; assigning a type to an ordered pair therefore fixes it for the
reversed pair.

Choose \(g_i:A\to A\) with \(\mathfrak A\entails\beta_i(a,g_i(a))\) for all
\(a\) and \(i\leq m\); then \(g_i(a)\neq a\).  Let \(K\) be the kings, \(P\)
the realized \(1\)-types, \(Q\) the types of kings, so \(\lvert Q\rvert=
\lvert K\rvert\), and
\(C\defeq K\cup\{g_i(k):k\in K,\ i\leq m\}\), the royal court.  Let
\(t_1,\ldots,t_n\) list \(P\setminus Q\), and take new elements
\(d_{ij},e_{ij},f_{ij}\) for \(i\leq m\), \(j\leq n\), forming the sets
\(D,E,F\).  The universe of \(\mathfrak B\) is \(C\cup D\cup E\cup F\); it is
nonempty, because \(P\neq\emptyset\), and if \(n=0\) every realized type
belongs to a king.

\emph{One-types.}  Each \(c\in C\) keeps its type \(t_c\) in \(\mathfrak A\),
and \(d_{ij},e_{ij},f_{ij}\) receive \(t_j\).  Distinct pairs inside \(C\)
keep their \(2\)-types in \(\mathfrak A\).  The kings of \(\mathfrak A\) are
kings of \(\mathfrak B\): an element of \(C\setminus K\) is not a king of
\(\mathfrak A\), so its type is in \(P\setminus Q\), and the new elements have
types in \(P\setminus Q\).

\emph{Witnesses and two-types.}  (1) For \(b\in K\), each \(g_i(b)\in C\) and
the pair already carries the type of \((b,g_i(b))\), so
\(\beta_i(b,g_i(b))\) holds in \(\mathfrak B\).  (2) For \(b\in C\setminus
K\) and \(i\leq m\): if \(g_i(b)\in C\), argue as in (1).  Otherwise
\(g_i(b)\) is not a king, its type is some \(t_j\), and we give the pair
\((b,d_{ij})\) the type of \((b,g_i(b))\); \(d_{ij}\) is the witness.  (3) For
\(b=d_{ij}\), fix \(a\in A\) of type \(t_j\).  For \(i'\leq m\): if
\(g_{i'}(a)\) is a king \(k\), give \((b,k)\) the type of \((a,k)\) and let
\(k\) be the witness; otherwise \(g_{i'}(a)\) has some type \(t_l\), give
\((b,e_{i'l})\) the type of \((a,g_{i'}(a))\) and let \(e_{i'l}\) be the
witness.  (4) Repeat (3) for \(b\in E\) with witnesses in \(K\cup F\), and for
\(b\in F\) with witnesses in \(K\cup D\).  In every case the assigned
\(2\)-type restricts to the \(1\)-types already given to its two members, and
contains \(x\neq y\).

No pair receives two different types.  Step (2) types only pairs in
\((C\setminus K)\times D\), with distinct \(i\) giving distinct second
members.  Step (3) types only pairs in \(D\times K\) and \(D\times E\); two
indices \(i'\) giving the same king use the same \((a,k)\) and so the same
type, and distinct \(i'\) give distinct elements of \(E\).  The \(E\)-step
types only \(E\times K\) and \(E\times F\), and the \(F\)-step only \(F\times
K\) and \(F\times D\).  These sets of unordered pairs are pairwise disjoint,
and no step types a pair inside \(C\).  This is where three banks are needed:
with two banks the pairs between \(D\) and \(E\) would be typed from both
sides.

(5) Every remaining pair \(\{b,b'\}\) of distinct elements receives the type
of a pair \((a,a')\) of distinct elements of \(\mathfrak A\) with
\(t_a=t_b\), \(t_{a'}=t_{b'}\).  Such a pair exists: if \(t_b\neq t_{b'}\) any
realizers are distinct; if \(t_b=t_{b'}\), this type is not a king type,
since a king type has only one element in \(\mathfrak B\), so it has two
realizers in \(\mathfrak A\).

\emph{Verification.}  Every element of \(\mathfrak B\) has a \(1\)-type, and
every pair of distinct elements a \(2\)-type, realized in \(\mathfrak A\).
Since \(\mathfrak A\entails\alpha(a,a')\) for all \(a,a'\), including
\(a=a'\), we get \(\mathfrak B\entails\forall x\forall y\alpha\).  Steps
(1)--(4) give each element a witness for each \(\beta_i\), carrying a
\(2\)-type in which \(\beta_i\) holds; so \(\mathfrak B\entails\forall
x\exists y\beta_i\) for every \(i\).

\emph{Counting.}  \(\lvert C\rvert\leq(m+1)\lvert K\rvert\) and
\(\lvert D\cup E\cup F\rvert=3m(\lvert P\rvert-\lvert K\rvert)\).  Since
\(m+1\leq3m\), \(\lvert B\rvert\leq3m\lvert P\rvert\).  Each conjunct uses a
symbol, so \(m\leq s\).  A \(1\)-type is fixed by one truth value per relation
symbol, namely \(Rx\), \(Rxx\) or the zero-ary \(R\), so
\(\lvert P\rvert\leq2^r\).  Hence \(\lvert B\rvert\leq3s2^r\)
\cite[proof of thm.~4.3]{GradelKolaitisVardi1997}.
\end{proof}
\end{annexitem}

\begin{annexitem}{annex:p02-w02-p02-035-fo2-emp}{thm:p02-w02-p02-035-fo2-emp}{Exponential model property for \(\mathrm{FO}^2\)}
\begin{proof}
Let \(\varphi\) be satisfiable of size \(s\), and let \(\varphi^{*}\) be its
Scott form from \cref{thm:p02-w02-p02-035-scott}, of size \(s^{*}=O(s\log s)\)
with \(r^{*}=O(s)\) relation symbols.  If \(\varphi^{*}\) has a one-element
model, so does \(\varphi\).  Otherwise \(\varphi^{*}\) has a model with at
least two elements, and \(m\geq1\), since a sentence of the form
\(\forall x\forall y\alpha\) with a model has a one-element model.  Let
\(\theta\) be \(\varphi^{*}\) with each \(\beta_i\) in the \(x\neq y\) form;
its size is \(O(s^{*})\), and it is true in that model.  By
\cref{thm:p02-w02-p02-035-scott-model}, \(\theta\) has a model \(\mathfrak B\)
with at most \(3\cdot O(s\log s)\cdot2^{O(s)}\) elements.  Since \(m\geq1\)
and each new \(\beta_i\) entails \(x\neq y\), \(\mathfrak B\) has at least two
elements, so \(\mathfrak B\entails\varphi^{*}\), and
\cref{thm:p02-w02-p02-035-scott} gives a model of \(\varphi\) on the same
universe.  The bound is at most \(2^{cs}\) for a suitable \(c\) and all
\(s\geq1\).  This is the source's Theorem~5.1
\cite[thm.~5.1]{GradelKolaitisVardi1997}.
\end{proof}
\end{annexitem}

\begin{annexitem}{annex:p02-w02-p02-035-fo2-complexity}{thm:p02-w02-p02-035-fo2-complexity}{\(\mathrm{FO}^2\) satisfiability is NEXPTIME-complete}
\begin{proof}
\emph{Scott-class input.}  Let \(\theta\) have size \(s\) and \(r\) relation
symbols, so \(r=O(s/\log s)\) by the counting fact.  Step 2 inspects
\(2^r\) one-element structures, each in time \(O(s)\).  If \(m=0\), a model
exists iff a one-element model exists, as in the previous proof.  Otherwise,
the modified sentence \(\theta'\) of Step 3 has size \(s'=O(s)\), and if it is
satisfiable it has a model with at most \(N=3s'2^r=2^{O(s/\log s)}\) elements.
Such a structure is a string of \(O(rN^2)\) truth values.  Step 4 evaluates
\(\alpha\) at \(N^2\) pairs and, for each of \(N\) elements and \(m\leq s\)
indices, tries at most \(N\) witnesses; each evaluation reads \(O(s)\) table
entries.  The total is \(2^{O(s/\log s)}\).

The procedure is correct: an accepted structure in Step 2 is a model of
\(\theta\); one accepted in Step 4 is a model of \(\theta'\) with at least two
elements, hence of \(\theta\).  If \(\theta\) has a model, either it has a
one-element model, found in Step 2, or a model with at least two elements, in
which \(\theta'\) holds, so
\cref{thm:p02-w02-p02-035-scott-model} provides a model of \(\theta'\) within
the guessed size.

\emph{\(\mathrm{FO}^2\) input.}  Step 1 is polynomial, and its output has size
\(s^{*}=O(s\log s)\) with \(r=O(s)\) symbols; satisfiability is preserved by
\cref{thm:p02-w02-p02-035-scott}.  The same analysis with
\(N=3s^{*}2^r=2^{O(s)}\) gives time \(2^{O(s)}\)
\cite[thm.~5.3]{GradelKolaitisVardi1997}.

\emph{Hardness.}  The sentence \(\psi_w\) of
\cref{def:p02-w02-p02-035-monadic-tiling-sentence}, renamed as in the proof of
\cref{thm:p02-w02-p02-035-monadic-hard}, is an equality-free
\(\mathrm{FO}^2\)-sentence of the form \(\forall x\exists y\alpha\land\forall
x\forall y\beta\) with only unary predicates.  So for
\(L\in\mathsf{NEXPTIME}\) the map \(x\mapsto\psi_{g(x)}\) of that proof
reduces \(L\) in polynomial time to satisfiability for this subclass.  The
subclass lies in \(\mathrm{FO}^2\) with and without equality, and the upper
bound covers \(\mathrm{FO}^2\) with equality, so both versions are
\(\mathsf{NEXPTIME}\)-complete.  By \cref{thm:p02-w02-p02-035-fo2-emp},
finite satisfiability is the same problem.
\end{proof}
\end{annexitem}

\begin{annexitem}{annex:p02-w02-p02-035-bsr-grounding}{thm:p02-w02-p02-035-bsr-grounding}{Correctness of grounding}
\begin{proof}
\emph{Model to valuation.}  Let \(\mathfrak A\entails\varphi\).  As in the
small-model argument of \cref{sec:p09-presentation-n-c-005}
\cite[prop.~6.2.17]{BorgerGradelGurevich1997}, choose
\(a_1,\ldots,a_p\) with \(\mathfrak A\entails\forall\bar y\theta[a_1,\ldots,
a_p]\); if \(p=0\) choose any \(a_1\).  Send \(c_i\) to \(a_i\) for \(i\leq q\).
Let \(c_i\equiv c_j\) iff \(a_i=a_j\); this is an equality pattern.  Give each
ground relation atom of \(G_\equiv\) the truth value of the corresponding
atom in \(\mathfrak A\); this is consistent because the least-indexed
representative of a class denotes the same element as every member.  Each
instance \(\theta(c_1,\ldots,c_p,c_{j_1},\ldots,c_{j_m})\) is true in
\(\mathfrak A\) under \(y_l\mapsto a_{j_l}\).  Replacing constants by
representatives and equality atoms by their truth values in \(\mathfrak A\),
which are those prescribed by \(\equiv\), does not change truth, so the
valuation satisfies \(G_\equiv\).  The induced substructure on
\(\{a_1,\ldots,a_q\}\) is itself a model with at most \(\max(1,p)\) elements.

\emph{Valuation to model.}  Let \(\equiv\) and a valuation \(v\) satisfy
\(G_\equiv\).  Let the universe \(B\) be the set of representatives, the
least-indexed member of each class.  For a relation symbol \(R\) put
\(R^{\mathfrak B}(\bar e)\) iff \(v(R\bar e)=1\) for representatives
\(\bar e\).  Interpret \(c_i\) as its representative \([c_i]\).  Then every
ground atom over \(c_1,\ldots,c_q\) has in \(\mathfrak B\) the value that
\(G_\equiv\) assigns to its replacement: relation atoms by the definition of
\(R^{\mathfrak B}\), and \(c_i=c_j\) because \([c_i]=[c_j]\) iff
\(c_i\equiv c_j\).  Every element of \(B\) is some \([c_j]\) with \(j\leq q\),
so every assignment of the \(y_l\) is \(y_l\mapsto[c_{j_l}]\) for some
instance in \(G(\varphi)\), which is true.  Hence
\(\mathfrak B\entails\forall\bar y\theta[[c_1],\ldots,[c_p]]\), and
\(\mathfrak B\entails\varphi\), with \(\lvert B\rvert\leq q=\max(1,p)\).
Without equality the identity pattern makes \(B=\{c_1,\ldots,c_q\}\).
\end{proof}
\end{annexitem}

\begin{annexitem}{annex:p02-w02-p02-035-bsr-complexity}{thm:p02-w02-p02-035-bsr-complexity}{Bernays--Sch\"onfinkel satisfiability is NEXPTIME-complete}
\begin{proof}
\emph{Upper bound.}  Let \(\varphi\) have length \(n\).  Then \(q\leq n\), and
the \(m\) universal variables have distinct names, so \(m=O(n/\log n)\) by the
counting fact.  The grounding has \(q^m\leq2^{m\log_2 n}=2^{O(n)}\) instances,
each of length \(O(n\log n)\) once constants are named in binary.  A pattern
is \(q^2\leq n^2\) bits.  Guessing a pattern and a valuation of the at most
\(2^{O(n)}\) ground atoms, and evaluating every instance, takes time
\(2^{O(n)}\).  Correctness is \cref{thm:p02-w02-p02-035-bsr-grounding}.

\emph{Lower bound.}  Fix \(\mathcal D=(D,H,V)\) and \(w\in D^n\), \(n\geq1\).
Use a unary symbol \(O\) and, for \(d\in D\), a \(2n\)-ary symbol \(P_d\).
Let \(\bar x,\bar y,\bar x',\bar y'\) be \(n\)-tuples of distinct variables.
For \(i<2^n\) write \(\bar u^{(i)}\) for the \(n\)-tuple whose \(l\)th entry is
\(u_1\) if bit \(l\) of \(i\) is \(1\) and \(u_0\) otherwise.  Let
\[
 \mathrm{Hor}\defeq\textstyle\bigwedge_{l<n}(Oy'_l\leftrightarrow Oy_l)\land
  \bigwedge_{l<n}\bigl(Ox'_l\leftrightarrow
   (Ox_l\oplus\bigwedge_{k<l}Ox_k)\bigr),
\]
and let \(\mathrm{Ver}\) be the same with the roles of the \(x\)- and
\(y\)-tuples exchanged.  Put
\[
 \chi_w\defeq\exists u_0\exists u_1\forall\bar x\forall\bar y\forall\bar x'
 \forall\bar y'\bigl(Ou_1\land\neg Ou_0\land\Gamma\bigr),
\]
where \(\Gamma\) states: exactly one \(P_d\bar x\bar y\) holds;
\(\mathrm{Hor}\to\bigvee_{(d,d')\in H}(P_d\bar x\bar y\land P_{d'}\bar x'\bar
y')\); \(\mathrm{Ver}\to\bigvee_{(d,d')\in V}(P_d\bar x\bar y\land
P_{d'}\bar x'\bar y')\); and \(P_{w_i}\bar u^{(i)}\bar u^{(0)}\) for \(i<n\).
This is a relational \(\exists\exists\forall^{*}\) sentence without equality
of length \(O(n^2\log n)\), computable from \(w\) in polynomial time.

If \(\tau\) witnesses \(w\in\mathrm{DOMINO}(\mathcal D,2^n)\), take universe
\(\{0,1\}\), \(O=\{1\}\), \(u_0\mapsto0\), \(u_1\mapsto1\), and let
\(P_d\) hold of \((\bar a,\bar b)\in\{0,1\}^{2n}\) iff \(\tau\) maps the point
with bit strings \(\bar a,\bar b\) to \(d\).  Every \(2n\)-tuple codes one
point, so exactly one \(P_d\) holds.  By the increment rule,
\(\mathrm{Hor}\) holds iff the second point is the right neighbour of the
first, and \(\mathrm{Ver}\) iff it is the upper neighbour; the tiling
conditions then give the implications, and \(\tau(i,0)=w_i\) gives the last
clauses.

Conversely, let \(\mathfrak A\entails\chi_w\) with witnesses \(e_0,e_1\) for
\(u_0,u_1\), so \(Oe_1\) and \(\neg Oe_0\).  For a point \(z\) let
\(\mathrm{tup}(z)\in\{e_0,e_1\}^{2n}\) spell its bits, and let \(\tau(z)\) be the
unique \(d\) with \(P_d\,\mathrm{tup}(z)\).  The truth values of \(O\) on
\(\mathrm{tup}(z)\) are the bits of \(z\), so instantiating the universal block
at \(\mathrm{tup}(z),\mathrm{tup}(z')\), with \(z'\) the right neighbour of
\(z\), makes \(\mathrm{Hor}\) true and yields \((\tau(z),\tau(z'))\in H\); the
vertical case is the same, and the last clauses give \(\tau(i,0)=w_i\).  So
\(\chi_w\) is satisfiable iff \(w\in\mathrm{DOMINO}(\mathcal D,2^n)\).
Composing with \cref{thm:p02-w02-p02-035-domino-hard} reduces every
\(L\in\mathsf{NEXPTIME}\) to this subclass in polynomial time.  The subclass
lies in the class with and without equality, so both are
\(\mathsf{NEXPTIME}\)-complete.  The negation of a \(\exists^{*}\forall^{*}\)
sentence is equivalent, by dualizing each quantifier in place, to a
\(\forall^{*}\exists^{*}\) sentence, so validity of the latter is
\(\mathsf{coNEXPTIME}\)-complete.
\end{proof}
\end{annexitem}

\begin{annexitem}{annex:p02-w02-p02-035-gf-constants}{thm:p02-w02-p02-035-gf-constants}{Constants can be removed}
\begin{proof}
If \(k=0\) put \(\psi'\defeq\psi\).  Otherwise read \(c_1,\ldots,c_k\) as
fresh distinct variables that occur nowhere in \(\psi\).  For every
subformula \(\theta\) of \(\psi\) let \(\theta^{*}\) replace each relation
atom \(R(t_1,\ldots,t_m)\), whose terms are variables or constants, by
\(R^{*}(c_1,\ldots,c_k,t_1,\ldots,t_m)\), and leave equalities unchanged
apart from reading constants as variables.

\emph{Guardedness.}  By induction on \(\theta\),
\(\theta^{*}\in\mathrm{GF}\) and
\(\mathrm{free}(\theta^{*})\subseteq\mathrm{free}(\theta)\cup\{\bar c\}\).
Atoms and Boolean steps are immediate.  For \(\theta=(\exists\bar y\,.\,\alpha)\chi\)
the guard \(\alpha^{*}=R^{*}(\bar c,\ldots)\) is a relation atom with
\(\mathrm{free}(\alpha^{*})=\mathrm{free}(\alpha)\cup\{\bar c\}\supseteq
\mathrm{free}(\chi)\cup\{\bar c\}\supseteq\mathrm{free}(\chi^{*})\); as
\(\bar y\) contains no \(c_j\), \((\exists\bar y\,.\,\alpha^{*})\chi^{*}\) is
built by rule (3).  The universal case is the same.  Since \(\psi\) is a
sentence, \(\mathrm{free}(\psi^{*})\subseteq\{\bar c\}=\mathrm{free}(Z\bar c)\),
so \(\psi'=(\exists\bar c\,.\,Z\bar c)\psi^{*}\) is a guarded sentence.  Each
atom grows by \(k\leq n\) symbols, so \(\lvert\psi'\rvert=O(n^2)\) and the
translation runs in polynomial time.

\emph{Satisfiability.}  Say that a structure \(\mathfrak A\) for \(\psi\) and
a structure \(\mathfrak B\) for \(\psi'\) with the same universe \emph{match
at} \(\bar e\in A^k\) if \(c_j^{\mathfrak A}=e_j\) for all \(j\) and
\(R^{\mathfrak A}=\{\bar u:(\bar e,\bar u)\in R^{*\mathfrak B}\}\) for every
\(R\).  If they match at \(\bar e\), then for every subformula \(\theta\) and
every assignment \(s\) of the variables of \(\psi\),
\(\mathfrak A\entails\theta[s]\) iff
\(\mathfrak B\entails\theta^{*}[s,\bar c\mapsto\bar e]\): for an atom both
sides say that the values of the terms lie in \(R^{\mathfrak A}\), Boolean
steps are immediate, and quantifiers range over the same universe while the
\(c_j\) are never bound.  Given \(\mathfrak A\entails\psi\), let
\(\mathfrak B\) have \(Z^{\mathfrak B}=\{\bar c^{\mathfrak A}\}\) and
\(R^{*\mathfrak B}=A^k\times R^{\mathfrak A}\); they match at
\(\bar c^{\mathfrak A}\), so \(\mathfrak B\entails\psi'\).  Given
\(\mathfrak B\entails\psi'\), choose \(\bar e\) with \(Z\bar e\) and
\(\mathfrak B\entails\psi^{*}[\bar c\mapsto\bar e]\), and let \(\mathfrak A\)
be the structure that matches \(\mathfrak B\) at \(\bar e\); then
\(\mathfrak A\entails\psi\).  Both structures have the same universe.
\end{proof}
\end{annexitem}

\begin{annexitem}{annex:p02-w02-p02-035-gf-normal-form}{thm:p02-w02-p02-035-gf-normal-form}{Normal form}
\begin{proof}
First rewrite every \((\forall\bar y\,.\,\alpha)\theta\) as
\(\neg(\exists\bar y\,.\,\alpha)\neg\theta\); this is guarded and equivalent.
Every formula contains an atom, and a constant-free sentence without
quantifiers can contain only \(0\)-ary relation atoms; since a guard
contains a variable, \(\varphi\) contains at least one relation atom.

\emph{Step 1: a common free variable.}  Choose a variable \(z\) not in
\(\varphi\), and replace every relation atom \(R(\bar u)\) by
\(R'(\bar u,z)\), where \(R'\) is fresh of arity \(\mathrm{ar}(R)+1\).  Call
the result \(\varphi'(z)\).  Guards stay relation atoms, and
\(\mathrm{free}(\theta')\subseteq\mathrm{free}(\theta)\cup\{z\}\subseteq
\mathrm{free}(\alpha)\cup\{z\}=\mathrm{free}(\alpha')\), so \(\varphi'\) is
guarded.  Every quantified subformula \((\exists\bar y\,.\,\alpha')\theta'\)
of \(\varphi'\) has \(z\) free, because its guard contains \(z\) and \(z\) is
never bound.  For a structure \(\mathfrak A\) and \(e\in A\), let
\(\mathfrak A'\) interpret \(R'\) as \(R^{\mathfrak A}\times A\).  Since
\(z\) is never bound, induction on formulas gives
\(\mathfrak A\entails\theta[s]\) iff \(\mathfrak A'\entails\theta'[s,z\mapsto e]\).
Conversely, if \(\mathfrak B\entails\varphi'[e]\), the structure with
\(R^{\mathfrak A}=\{\bar u:(\bar u,e)\in R'^{\mathfrak B}\}\) satisfies the
same equivalence, hence \(\varphi\).  So \(\varphi\) has a model with
universe \(A\) iff \(\varphi'\) is satisfied in a structure with universe
\(A\) at some \(e\).

\emph{Step 2: naming quantified subformulas.}  Maintain a guarded formula
\(\varphi_{\mathrm{cur}}(z)\) and a conjunction \(\eta\) of sentences of the
normal-form shape, starting with \(\varphi'\) and the empty conjunction.
While \(\varphi_{\mathrm{cur}}\) contains a quantifier, choose an occurrence
of a subformula \(\chi=(\exists\bar y\,.\,\alpha)\theta\) with \(\theta\)
quantifier-free, list \(\mathrm{free}(\chi)=\mathrm{free}(\alpha)\setminus\bar y\)
as \(\bar z\) (so \(z\in\bar z\)), take a fresh \(R_\chi\) of arity
\(\lvert\bar z\rvert\), replace the occurrence by \(R_\chi\bar z\), and add to
\(\eta\) the two sentences
\[
 \eta_\chi\defeq(\forall\bar z\,.\,R_\chi\bar z)\exists\bar y\,(\alpha\land\theta),
 \qquad
 \eta'_\chi\defeq(\forall\bar z\bar y\,.\,\alpha)(\theta\to R_\chi\bar z).
\]
Both have the normal-form shape: in \(\eta_\chi\) the guard has exactly the
variables \(\bar z\); in \(\eta'_\chi\) the guard \(\alpha\) has exactly the
variables \(\bar z\bar y\) and the existential tuple is empty.  The guard of
every remaining quantified subformula is an original guard of
\(\varphi'\), so it contains \(z\); and the replacement does not change free
variables, so \(\varphi_{\mathrm{cur}}\) stays guarded.  Each step removes
one quantifier, so the loop ends with a quantifier-free \(\varphi''(z)\).

Each step preserves satisfiability over a fixed universe.  If
\(\mathfrak B\entails\varphi_{\mathrm{cur}}[e]\land\eta\), expand
\(\mathfrak B\) by \(R_\chi^{\mathfrak B}\defeq\{\bar c:\mathfrak B\entails\chi[\bar c]\}\);
then \(\eta_\chi\) and \(\eta'_\chi\) hold and the new formula has the same
truth value as the old one at every assignment.  Conversely, in any
structure satisfying \(\eta_\chi\land\eta'_\chi\), \(R_\chi\bar c\) implies
\(\chi[\bar c]\) by \(\eta_\chi\), and \(\chi[\bar c]\) with a witness
\(\bar d\) gives \(\alpha[\bar c,\bar d]\) and \(\theta[\bar c,\bar d]\),
hence \(R_\chi\bar c\) by \(\eta'_\chi\).  So \(R_\chi\) defines \(\chi\),
and the new formula agrees with the old one.

\emph{Step 3.}  Put \(\psi\defeq\exists z\,Cz\land(\forall z\,.\,Cz)\varphi''(z)\land\eta\)
with a fresh unary \(C\); the middle conjunct has the normal-form shape with
empty existential tuple.  If \(\varphi''\land\eta\) holds at \(e\), interpret
\(C\) as \(\{e\}\); conversely pick \(e\in C\).  Renaming bound variables
puts every conjunct into the indexing of
\cref{def:p02-w02-p02-035-gf-normal}.

\emph{Bounds.}  The variables of \(\psi\) are those of \(\varphi\) and \(z\).
The symbols \(R'\) have arity at most \(a+1\); \(R_\chi\) has arity
\(\lvert\bar z\rvert\), at most the number of distinct variables of a guard
\(R'(\ldots)\), hence at most \(a+1\); \(C\) is unary.  There are at most
\(n\) steps.  Replacing \(\chi\) by \(R_\chi\bar z\) does not lengthen
\(\varphi_{\mathrm{cur}}\) when symbol names are counted as single symbols,
because \(\bar z\) is a subtuple of the variables of \(\alpha\).  Each step
therefore adds \(O(n)\) symbols, and \(\psi\) has \(O(n^2)\) symbols, each
written with \(O(\log n)\) bits.  All steps are polynomial-time.
\end{proof}
\end{annexitem}

\begin{annexitem}{annex:p02-w02-p02-035-gf-tree}{thm:p02-w02-p02-035-gf-tree}{Witness criterion and tree-model property}
\begin{proof}
Write \(\mathrm{At}_k\) for \(\mathrm{At}_k(\psi)\) and \(\tau\) for the
vocabulary of \(\psi\).  For a map \(\rho\) between index sets write
\(\rho\lambda\) for the atom obtained by replacing each \(x_p\) by
\(x_{\rho(p)}\).  Five facts are used.

(F1) If \(\lambda\in\mathrm{At}_k\) and \(\rho:\{1,\ldots,k\}\to\{1,\ldots,k'\}\)
with \(k'\leq m\), then \(\rho\lambda\in\mathrm{At}_{k'}\); and an atom of
\(\mathrm{At}_{k'}\) whose indices are all at most \(k\) lies in
\(\mathrm{At}_k\).  Both follow by composing or restricting the map \(h\) in
\cref{def:p02-w02-p02-035-gf-types}.

(F2) If \(\bar a\) realizes a \(k'\)-type \(t\), then
\(\bar a\circ\rho=(a_{\rho(1)},\ldots,a_{\rho(k)})\) realizes
\(t^\rho\defeq\{\pm\lambda:\lambda\in\mathrm{At}_k,\ \pm\rho\lambda\in t\}\),
where \(\pm\) is the same sign on both sides.  So the reduction along
\(\rho\) exists and is unique.  If \(k\leq k'\) and \(\rho\) is the inclusion,
\(t^\rho=t\cap(\mathrm{At}_k\cup\neg\mathrm{At}_k)\); hence \(t\) extends
\(s\) iff \(s\) is this restriction of \(t\).

(F3) For a type \(t\), let \(p\approx_t q\) iff \((x_p=x_q)\in t\).  If
\(\bar a\) realizes \(t\), then \(p\approx_t q\) iff \(a_p=a_q\).  Hence, if
\(\rho(p)\approx_t p\) for all \(p\) and both \(\lambda\) and \(\rho\lambda\)
lie in \(\mathrm{At}_k\), then \(\lambda\in t\) iff \(\rho\lambda\in t\).

(F4) A quantifier-free formula all of whose atoms lie in \(\mathrm{At}_k\) is
true of \(\bar a\) iff it is true under the valuation that
\(\mathrm{atp}^{\psi}(\bar a)\) gives its atoms.  The atoms of
\(C\bar x\), \(\alpha_i\) and \(\psi_i\), indexed as in
\cref{def:p02-w02-p02-035-gf-normal}, lie in \(\mathrm{At}_r\) and
\(\mathrm{At}_{k_i+\ell_i}\) by definition.

(F5) Call \(R(b_1,\ldots,b_r)\) \emph{relevant} if some atom
\(R(z_1,\ldots,z_r)\) of \(\psi\) satisfies \(b_p=b_q\) whenever
\(z_p=z_q\).  If \(\bar e\) lists a set containing \(b_1,\ldots,b_r\) and
\(\bar b=\bar e\circ\bar p\), then \(R(x_{p_1},\ldots,x_{p_r})\in\mathrm{At}_k\)
implies that \(R(\bar b)\) is relevant.  Conversely, if \(R(\bar b)\) is
relevant, choosing one index for each element gives \(\bar p\) with
\(p_c=p_d\) iff \(b_c=b_d\), and then \(R(x_{\bar p})\in\mathrm{At}_k\).

\emph{(a)\(\Rightarrow\)(b).}  Let \(\mathfrak A\entails\psi\).  Let \(P\) be
the set of \(k\)-types, \(1\leq k\leq m\), realized in \(\mathfrak A\); it
is finite and, by (F2), closed under reductions.  Choose \(\bar b\) with
\(\mathfrak A\entails C\bar b\) and put
\(t_0\defeq\mathrm{atp}^{\psi}(\bar b)\).  For each \(s\in P^{(k_i)}\) fix a
realization \(\bar a_s\).  If \(\alpha_i\in s\), then
\(\mathfrak A\entails\alpha_i(\bar a_s)\) by (F4), so some \(\bar d\) has
\(\mathfrak A\entails\psi_i(\bar a_s,\bar d)\); otherwise let \(\bar d\) repeat
the first entry of \(\bar a_s\).  Put
\(\mathrm{ext}_i(s)\defeq\mathrm{atp}^{\psi}(\bar a_s\bar d)\).  It lies in
\(P\) because \(k_i+\ell_i\leq m\), extends \(s\) by (F2), and makes
\(\alpha_i\to\psi_i\) true by (F4).

\emph{(b)\(\Rightarrow\)(c): the unravelling.}  Fix a witness
\((P,t_0,(\mathrm{ext}_i))\).  We build a rooted tree \(T\) and, for each node
\(v\), a type \(t_v\in P^{(k_v)}\) and a surjection \(e_v\) from
\(\{1,\ldots,k_v\}\) onto a set \(F(v)\) with
\(e_v(p)=e_v(q)\) iff \(p\approx_{t_v}q\).  Each element is \emph{new} at
exactly one node.  At the root \(\lambda\), \(t_\lambda\defeq t_0\), and
\(F(\lambda)\) has one fresh element for each \(\approx_{t_0}\)-class, all new
at \(\lambda\).  Given \(v\), for each \(i\in I\) and each
\(\sigma:\{1,\ldots,k_i\}\to\{1,\ldots,k_v\}\) such that (1) some
\(e_v(\sigma(p))\) is new at \(v\), (2) \(\alpha_i\in s\defeq t_v^\sigma\),
and (3) \(t\defeq\mathrm{ext}_i(s)\) extends \(s\) by \(j\geq1\) new
elements, create a child \(w\) with \(t_w\defeq t\).  For \(p\leq k_i\) put
\(e_w(p)\defeq e_v(\sigma(p))\).  For \(q>k_i\) put \(e_w(q)\defeq e_v(\sigma(p))\)
if \(q\approx_t p\) for some \(p\leq k_i\), and otherwise let \(e_w(q)\) be a
fresh element for the class of \(q\); the fresh elements are new at \(w\).
This is well defined and satisfies \(e_w(p)=e_w(q)\) iff \(p\approx_t q\):
for \(p,p'\leq k_i\), \(p\approx_t p'\) iff \((x_p=x_{p'})\in s\) by (F2)
iff \((x_{\sigma p}=x_{\sigma p'})\in t_v\) iff
\(e_v(\sigma p)=e_v(\sigma p')\); the remaining cases follow by transitivity
of \(\approx_t\) and the choice of one fresh element per class.  Every node
has at most \(\lvert I\rvert m^m\) children, and \(\lvert F(v)\rvert\leq m\).

\emph{Connectivity.}  An element of \(F(w)\) that is not new at \(w\) lies in
\(F(v)\) for the parent \(v\).  So the nodes containing an element \(b\) are
the node where \(b\) is new together with descendants of it, each of whose
parents also contains \(b\); this set is connected.  This is condition (ii)
of a \(k\)-tree.

\emph{The structure.}  Let \(B\defeq\bigcup_vF(v)\), and for \(R\in\tau\) let
\(R^{\mathfrak B}\) be the set of relevant tuples \(\bar b\) such that, for
some node \(v\) and some \(\bar p\), \(\bar b=e_v\circ\bar p\) and
\(R(x_{\bar p})\in t_v\cap\mathrm{At}_{k_v}\).

\emph{Claim 1.}  If \(\bar b=e_v\circ\bar p\) and
\(R(x_{\bar p})\in\mathrm{At}_{k_v}\), then \(\mathfrak B\entails R(\bar b)\)
iff \(R(x_{\bar p})\in t_v\).  Within one node, two index tuples
\(\bar p,\bar p'\) with \(e_v\circ\bar p=e_v\circ\bar p'\) satisfy
\(p_c\approx_{t_v}p'_c\), so by (F3) \(R(x_{\bar p})\in t_v\) iff
\(R(x_{\bar p'})\in t_v\).  For a node \(v\) and its child \(w\) built from
\(\sigma\), \(F(v)\cap F(w)=\{e_v(\sigma p):p\leq k_i\}\).  For a relevant
\(\bar b\) in this set, (F5) gives \(\bar q\) with entries at most \(k_i\),
\(e_w\circ\bar q=\bar b\) and \(R(x_{\bar q})\in\mathrm{At}_{k_i}\).  Then
\(R(x_{\bar q})\in t_w\) iff \(R(x_{\bar q})\in s\) by (F2), iff
\(R(x_{\sigma\bar q})\in t_v\) by definition of \(s=t_v^\sigma\), and
\(e_v\circ\sigma\circ\bar q=\bar b\) with \(R(x_{\sigma\bar q})\in\mathrm{At}_{k_v}\)
by (F1).  So adjacent nodes agree.  If \(\bar b\) lies in \(F(u)\) and in
\(F(v)\), then by connectivity every entry of \(\bar b\) lies in every set
\(F(v')\) on the path from \(u\) to \(v\), and agreement passes along the
path.  This proves Claim~1.

\emph{Claim 2.}  For every node \(v\), \(e_v\) realizes \(t_v\) in
\(\mathfrak B\).  Equalities hold by construction.  For
\(R(x_{\bar p})\in\mathrm{At}_{k_v}\) the tuple \(e_v\circ\bar p\) is relevant
by (F5), and Claim~1 applies.  By (F2), every \(k\)-tuple over \(F(v)\), with
\(k\leq m\), realizes a reduction of \(t_v\), which lies in \(P\).

By definition every tuple in some \(R^{\mathfrak B}\) lies in one set
\(F(v)\), which is condition (i).  So \(\mathfrak B\) is an \(m\)-tree of
bounded degree.

\emph{\(\mathfrak B\entails\psi\).}  Since \(C\bar x\in t_0=t_\lambda\),
Claim~2 gives \(\mathfrak B\entails C(e_\lambda)\).  Fix \(i\) and
\(\bar a\in B^{k_i}\) with \(\mathfrak B\entails\alpha_i(\bar a)\).  The
argument tuple of \(\alpha_i\) at \(\bar a\) lies in \(R^{\mathfrak B}\) for
the symbol \(R\) of \(\alpha_i\), so it lies in one set \(F(v)\); it lists
every entry of \(\bar a\).  The set of nodes containing all entries of
\(\bar a\) is an intersection of connected sets of a tree, hence connected
and nonempty, so it has a unique node \(v_0\) nearest the root.  Some entry
of \(\bar a\) is new at \(v_0\): otherwise all entries lie in the parent of
\(v_0\), and at the root every element is new.  Write
\(\bar a=e_{v_0}\circ\sigma\) and \(s\defeq t_{v_0}^\sigma\), the type of
\(\bar a\) by Claim~2 and (F2); then \(\alpha_i\in s\) by (F4).  Let
\(t\defeq\mathrm{ext}_i(s)\), so \(\psi_i\) is true in \(t\).  If \(t\)
extends \(s\) by \(j\geq1\) new elements, the child \(w\) built from
\(v_0,i,\sigma\) exists; \(e_w\) realizes \(t\) and agrees with \(\bar a\) on
\(1,\ldots,k_i\), so \(\bar d\defeq(e_w(k_i+1),\ldots,e_w(k_i+\ell_i))\)
satisfies \(\mathfrak B\entails\psi_i(\bar a,\bar d)\) by (F4).  If \(j=0\),
every \(\approx_t\)-class meets \(\{1,\ldots,k_i\}\); let \(\rho\) fix
\(1,\ldots,k_i\) and send each \(q>k_i\) to some \(p\leq k_i\) with
\(p\approx_tq\).  Then \(\bar a\circ\rho\) realizes \(s^\rho\), and
\(s^\rho=t\): for \(\lambda\in\mathrm{At}_{k_i+\ell_i}\), \(\lambda\in s^\rho\)
iff \(\rho\lambda\in s\) iff \(\rho\lambda\in t\) (as \(s\subseteq t\) and
\(\rho\lambda\in\mathrm{At}_{k_i}\)) iff \(\lambda\in t\) by (F3).  So
\(\bar d\defeq(a_{\rho(k_i+1)},\ldots,a_{\rho(k_i+\ell_i)})\) satisfies
\(\psi_i(\bar a,\bar d)\).  This covers \(\ell_i=0\) as well.  Hence
\(\mathfrak B\entails\psi\).

\emph{(c)\(\Rightarrow\)(a)} is immediate.  The last sentence of the theorem
follows from \cref{thm:p02-w02-p02-035-gf-normal-form}, which gives the
normal form at most one more variable.
\end{proof}
\end{annexitem}

\begin{annexitem}{annex:p02-w02-p02-035-gf-fmp}{thm:p02-w02-p02-035-gf-fmp}{Finite-model property}
\begin{proof}
Let \(\varphi\) be a satisfiable guarded sentence.  By
\cref{thm:p02-w02-p02-035-gf-constants,thm:p02-w02-p02-035-gf-normal-form}
there is a sentence \(\psi\) in normal form, over a finite relational
vocabulary \(\tau\), such that \(\varphi\) and \(\psi\) have models with the
same universes; it suffices to give \(\psi\) a finite model.  By
\cref{thm:p02-w02-p02-035-gf-tree}, \(\psi\) has a witness; let
\(\mathfrak B,T,F,t_v,e_v\) be the unravelling built from it in that proof,
and let \(\mathfrak F(v)\) be the substructure of \(\mathfrak B\) induced on
\(F(v)\).

Since \(\lvert F(v)\rvert\leq m\) and \(\tau\) is finite, the structures
\(\mathfrak F(v)\) fall into finitely many isomorphism types.  Choose one
node of each type and add all their predecessors; the result is a finite set
\(S\) of nodes, closed under taking parents, such that every \(\mathfrak F(w)\)
is isomorphic to some \(\mathfrak F(u)\) with \(u\in S\).  Let \(\mathfrak A\)
be the substructure of \(\mathfrak B\) induced on
\(A\defeq\bigcup_{v\in S}F(v)\), and let \(\mathfrak A^{+}\) be as in
\cref{def:p02-w02-p02-035-herwig-import}.  We show
\(\mathfrak A^{+}\entails\psi\).

The root lies in \(S\) and \(\mathfrak B\entails C(e_\lambda)\); as
\(\mathfrak A\) is induced in \(\mathfrak B\) and in \(\mathfrak A^{+}\),
\(\mathfrak A^{+}\entails\exists\bar x\,C\bar x\).

Fix \(i\) and \(\bar a\) with \(\mathfrak A^{+}\entails\alpha_i(\bar a)\).
The argument tuple of \(\alpha_i\) at \(\bar a\) satisfies a relation atom in
which all its positions are distinct variables, so it is live and lists every
entry of \(\bar a\).  By premise (ii) an automorphism \(g\) moves it into
\(A\); then \(\bar b\defeq g(\bar a)\) lies in \(A\), and
\(\mathfrak A^{+}\entails\alpha_i(\bar b)\), hence
\(\mathfrak A\entails\alpha_i(\bar b)\) and \(\mathfrak B\entails\alpha_i(\bar b)\).

Let \(v_0\) be the node nearest the root among the nodes containing all
entries of \(\bar b\), as in the tree proof.  We show \(v_0\in S\).  Each entry
lies in \(A\), so in some \(F(u)\) with \(u\in S\); the node where it is new
is a predecessor of \(u\) or \(u\) itself, so lies in \(S\).  These nodes,
one per entry, are all predecessors of any node containing \(\bar b\), so
they lie on one path from the root.  Let \(\beta\) be the deepest of them.
Every entry is new at a predecessor of \(\beta\) and present below \(\beta\),
so by connectivity it is present at \(\beta\); no proper predecessor of
\(\beta\) contains the entry that is new at \(\beta\).  So \(v_0=\beta\in S\).

As in the tree proof, \(\bar b=e_{v_0}\circ\sigma\), \(s=t_{v_0}^\sigma\)
contains \(\alpha_i\), and \(t=\mathrm{ext}_i(s)\) makes \(\psi_i\) true.

If \(t\) adds no new elements, the tree proof found
\(\bar d\subseteq F(v_0)\subseteq A\) with \(\mathfrak B\entails\psi_i(\bar b,\bar d)\).
Since \(\psi_i\) is quantifier-free and the substructures are induced,
\(\mathfrak A^{+}\entails\psi_i(\bar b,\bar d)\), and applying \(g^{-1}\) gives
\(\mathfrak A^{+}\entails\psi_i(\bar a,g^{-1}\bar d)\).

Otherwise the child \(w\) of \(v_0\) built from \(i,\sigma\) carries
\(\bar d\subseteq F(w)\) with \(\mathfrak F(w)\entails\psi_i(\bar b,\bar d)\).
Choose \(u\in S\) and an isomorphism \(\pi:\mathfrak F(w)\to\mathfrak F(u)\).
Let \(p\) be the restriction of \(\pi\) to the entries of \(\bar b\), which
lie in \(F(v_0)\cap F(w)\subseteq A\).  Then \(p\) is a partial isomorphism of
\(\mathfrak A\): it is injective, and for every relation symbol \(R\) and
tuple \(\bar c\) over its domain, \(\mathfrak A\entails R(\bar c)\) iff
\(\mathfrak F(w)\entails R(\bar c)\) iff \(\mathfrak F(u)\entails R(\pi\bar c)\)
iff \(\mathfrak A\entails R(\pi\bar c)\), all substructures being induced.
By premise (i) some automorphism \(f\) of \(\mathfrak A^{+}\) extends \(p\).
Now \(\mathfrak F(u)\entails\psi_i(\pi\bar b,\pi\bar d)\), so
\(\mathfrak A^{+}\entails\psi_i(f g\bar a,\pi\bar d)\), and applying
\((fg)^{-1}\) gives \(\mathfrak A^{+}\entails\psi_i(\bar a,(fg)^{-1}\pi\bar d)\).

In both cases \(\mathfrak A^{+}\entails\exists\bar y\,\psi_i(\bar a,\bar y)\).
So the finite structure \(\mathfrak A^{+}\) is a model of \(\psi\), and
\(\varphi\) has a model with the same finite universe.  Since every finite
model is a model, finite and unrestricted satisfiability coincide.
\end{proof}
\end{annexitem}

\begin{annexitem}{annex:p02-w02-p02-035-gf-upper}{thm:p02-w02-p02-035-gf-upper}{Correctness and running time of type elimination}
\begin{proof}
\emph{Step 2 lists exactly the types.}  A type satisfies both conditions by
(F3) of the tree proof.  Conversely, let \(t\) satisfy them, with
equivalence \(\approx\).  Let \(\mathfrak M\) have universe
\(\{1,\ldots,k\}/{\approx}\) and
\(R^{\mathfrak M}\defeq\{([p_1],\ldots,[p_r]):R(x_{\bar p})\in t\}\).  The
tuple \(([1],\ldots,[k])\) satisfies the equalities of \(t\) and its
positive atoms.  If \(\neg R(x_{\bar p})\in t\) but
\(([p_1],\ldots,[p_r])\in R^{\mathfrak M}\), then some \(R(x_{\bar q})\in t\)
has \(q_c\approx p_c\) for all \(c\), against the second condition.  So
\(t\) is realized.

\emph{Correctness.}  The sets \(P_j\) decrease and are finite, so the loop
stops.  Suppose the output is positive.  At the fixed point, every
reduction of a member of \(P^{*}\) lies in \(P^{*}\).  For \(s\in P^{*}\) a
\(k_i\)-type with \(\alpha_i\in s\), let \(\mathrm{ext}_i(s)\) be a type
\(t'\) found in Step~3.  For \(s\) without \(\alpha_i\), let
\(\mathrm{ext}_i(s)\defeq s^\rho\) with \(\rho\) fixing \(1,\ldots,k_i\) and
sending every \(q>k_i\) to \(1\); it is a reduction of \(s\), hence in
\(P^{*}\), extends \(s\) by (F2), and makes \(\alpha_i\to\psi_i\) true
because \(\neg\alpha_i\in s\).  With \(t_0\) from the output this is a
witness, so \(\psi\) is satisfiable by \cref{thm:p02-w02-p02-035-gf-tree},
and \(\varphi\) is satisfiable by
\cref{thm:p02-w02-p02-035-gf-constants,thm:p02-w02-p02-035-gf-normal-form}.
Conversely, let \(\mathfrak A\entails\psi\) and let \(Q\) be the set of types,
with at most \(m\) variables, realized in \(\mathfrak A\).  Then
\(Q\subseteq P_0\).  If \(Q\subseteq P_j\) and \(t\in Q\) is realized by
\(\bar a\), every reduction of \(t\) is realized and so lies in \(Q\); and
if \(\alpha_i\in t\), some \(\bar d\) has \(\mathfrak A\entails\psi_i(\bar a,\bar d)\),
so \(\mathrm{atp}^{\psi}(\bar a\bar d)\in Q\) extends \(t\) and makes
\(\psi_i\) true.  Hence \(Q\subseteq P_{j+1}\), so \(Q\subseteq P^{*}\), and
the type of a tuple in \(C^{\mathfrak A}\) lies in \(P^{*}\).

\emph{Running time.}  Let \(m\) be the number of variables of \(\psi\), \(q\)
the largest number of distinct variables in an atom of \(\psi\), and
\(K\defeq\lvert\psi\rvert m^{q}+m^2\).  Each atom of \(\psi\) yields at most
\(m^{q}\) members of \(\mathrm{At}_k\), so \(\lvert\mathrm{At}_k\rvert\leq K\)
and there are at most \(N\defeq m2^{K}\) types; each is written with a
number of bits polynomial in \(K\) and \(\lvert\psi\rvert\).  Step~2 takes
time \(N\) times a polynomial in \(K\).  Step~3 has at most \(N+1\) rounds.
In one round, for each of at most \(N\) types it computes at most
\(m^{m+1}\) reductions and tests each for membership among at most \(N\)
types, and for each of at most \(\lvert\psi\rvert\) indices \(i\) it tests at
most \(N\) candidates \(t'\).  Each test is polynomial in \(K\) and
\(\lvert\psi\rvert\).  The total time is therefore
\(N^{3}m^{m+1}\) times a polynomial in \(K\) and \(\lvert\psi\rvert\), that is
\(2^{O(K+m\log m)}\).

In general, removing \(k\leq n\) constants gives at most \(2n\) variables, and
the normal form adds one, so \(m\leq2n+1\), \(q\leq m\), and
\(\lvert\psi\rvert\) is polynomial in \(n\) by
\cref{thm:p02-w02-p02-035-gf-constants,thm:p02-w02-p02-035-gf-normal-form}.
Then \(K\leq\lvert\psi\rvert(2n+1)^{2n+1}+(2n+1)^2=2^{O(n\log n)}\), and the
running time is \(2^{2^{O(n\log n)}}\).  Since \(cn\log n\leq n^2\) for large
\(n\), this lies in \(\mathsf{DTIME}(2^{2^{n^2}})\) up to finitely many
inputs, so the problem is in \(2\mathsf{EXPTIME}\).

For constant-free input with at most \(k\) variables, no constants are
removed, \(m\leq k+1\) and \(q\leq k+1\), so \(K\leq\lvert\psi\rvert(k+1)^{k+1}+(k+1)^2\)
is polynomial in \(n\) and \(m\log m\) is constant.  For constant-free input
whose relation symbols have arity at most \(k\), the normal form has
arities at most \(k+1\), so \(q\leq k+1\), \(m\leq n+1\), \(K\) is polynomial
in \(n\), and \(m\log m=O(n\log n)\).  In both cases the time is
\(2^{n^{O(1)}}\), so the problems are in \(\mathsf{EXPTIME}\).
\end{proof}
\end{annexitem}

\begin{annexitem}{annex:p02-w02-p02-035-gf-bounded-constants}{thm:p02-w02-p02-035-gf-bounded-constants}{Bounded guarded sentences with constants}
\begin{proof}
Fix \(k\geq1\) and put \(q\defeq k+1\).  Let \(\varphi\) be a guarded sentence
of length \(n\) whose constants are among \(c_1,\ldots,c_p\), \(p\leq n\), and
which either has at most \(k\) variables or has only relation symbols of
arity at most \(k\).

\emph{Step 1: an unnamed element.}  Let \(\mathfrak A\) be a structure for
\(\varphi\), and let \(\mathfrak A'\) add to it one new element \(e\) that lies
in no relation, with the constants and relations of \(\mathfrak A\)
unchanged.  We show by induction that for every guarded formula \(\theta\)
and every assignment \(s\) into \(A\), \(\mathfrak A\entails\theta[s]\) iff
\(\mathfrak A'\entails\theta[s]\).  Atoms have the same value tuples and the
same relations in both structures, and Boolean steps are immediate.  Let
\(\theta=(\exists\bar y\,.\,\alpha)\chi\).  If \(\bar b\) is a tuple over
\(A'\) with \(\mathfrak A'\entails(\alpha\land\chi)[s,\bar y\mapsto\bar b]\),
then every variable of \(\bar y\) occurs in the relation atom \(\alpha\), so
every entry of \(\bar b\) is an entry of a tuple of
\(R^{\mathfrak A'}=R^{\mathfrak A}\subseteq A^{r}\), where \(R\) is the symbol of
\(\alpha\); hence \(\bar b\) lies over \(A\), and the induction hypothesis gives
\(\mathfrak A\entails(\alpha\land\chi)[s,\bar y\mapsto\bar b]\).  The
converse uses the same \(\bar b\).  For
\(\theta=(\forall\bar y\,.\,\alpha)\chi\), a tuple \(\bar b\) over \(A'\) with
\(\mathfrak A'\entails\alpha[s,\bar y\mapsto\bar b]\) lies over \(A\) by the
same argument, so both structures test \(\chi\) on the same tuples, and the
induction hypothesis applies to each.  Hence \(\varphi\) has a model iff it
has a model in which some element is the value of no constant.

\emph{Step 2: an equality pattern.}  For an equivalence relation \(\equiv\)
on \(\{1,\ldots,p\}\) let \(\varphi_\equiv\) replace every \(c_j\) by
\(c_{j'}\), where \(j'\) is the least member of the class of \(j\).  This is
a guarded sentence with the same variables and relation symbols.  If
\(\mathfrak A\entails\varphi\) and \(\equiv\) relates \(i\) and \(j\) exactly
when \(c_i^{\mathfrak A}=c_j^{\mathfrak A}\), then
\(\mathfrak A\entails\varphi_\equiv\), since each replaced constant has the
same value as its replacement, and the remaining constants have pairwise
distinct values.  Conversely, if \(\mathfrak B\entails\varphi_\equiv\),
interpreting each \(c_j\) as \(c_{j'}^{\mathfrak B}\) gives a model of
\(\varphi\) with the same universe.  Adding an element as in Step 1 keeps the
values of the constants distinct.  So \(\varphi\) is satisfiable iff for some
\(\equiv\) the sentence \(\varphi_\equiv\) has a model in which its constants
have pairwise distinct values and some element is the value of no constant.
Below, the constants of \(\varphi_\equiv\) are renamed \(c_1,\ldots,c_p\).

\emph{Step 3: flattening.}  Apply to \(\varphi_\equiv\) the rewriting of
universal guarded quantifiers and Steps 1 and 2 of the proof of
\cref{thm:p02-w02-p02-035-gf-normal-form}.  Those steps remain valid when
atoms contain constants.  In Step 1, \(R(\bar u)\) becomes \(R'(\bar u,z)\)
whatever terms \(\bar u\) contains; the guardedness check concerns free
variables, to which constants contribute nothing; and in the semantic
induction the value tuple of \(R(\bar u)\) in \(\mathfrak A\) followed by
\(e\) is the value tuple of \(R'(\bar u,z)\) in \(\mathfrak A'\) when both
structures interpret the constants alike.  In Step 2, \(\bar z\) lists the
free variables of \(\chi\), which are variables, and the constants of
\(\alpha\) and \(\theta\) stay in \(\eta_\chi\) and \(\eta'_\chi\); the argument
that \(R_\chi\) defines \(\chi\) evaluates formulas in one structure and does
not depend on which terms occur.  The expansions and reducts in that proof
leave the constants unchanged.  Put
\[
 \varphi^{\flat}\defeq(\exists z\,.\,Cz)Cz\land(\forall z\,.\,Cz)\varphi''(z)\land\eta
\]
with a fresh unary \(C\).  Then \(\varphi_\equiv\) has a model with universe
\(A\) and given values of the constants iff \(\varphi^{\flat}\) does: if
\(\varphi''\land\eta\) holds at \(e\), interpret \(C\) as \(\{e\}\); conversely
pick \(e\in C\).  Every conjunct of \(\varphi^{\flat}\) has one of three
shapes: \((\exists z\,.\,Cz)Cz\); \((\forall\bar x\,.\,\alpha)\theta\) with
\(\theta\) quantifier-free and \(\mathrm{free}(\alpha)=\{\bar x\}\), which covers
\((\forall z\,.\,Cz)\varphi''\) and every \(\eta'_\chi\); and
\((\forall\bar x\,.\,\alpha)(\exists\bar y\,.\,\beta)\theta\) with \(\theta\)
quantifier-free, \(\bar y\) nonempty, \(\mathrm{free}(\alpha)=\{\bar x\}\) and
\(\mathrm{free}(\beta)=\{\bar x,\bar y\}\), which covers every \(\eta_\chi\).
The variables of \(\varphi^{\flat}\) are those of \(\varphi\) and \(z\); the
symbols \(R'\) and \(R_\chi\) have arity at most one more than the largest
arity of \(\varphi\), as in that proof, and \(C\) is unary.  So in both
bounded cases every atom of \(\varphi^{\flat}\) has at most \(q\) distinct
variables, and \(\varphi^{\flat}\) has length polynomial in \(n\).

\emph{Step 4: the split.}  Let \(\ast\) be a new symbol.  For a relation symbol
\(R\) of arity \(r\) and a map \(\kappa:\{1,\ldots,r\}\to\{\ast,1,\ldots,p\}\)
with \(\kappa^{-1}(\ast)=\{p_1<\cdots<p_l\}\) nonempty, let \(R_\kappa\) be a
fresh relation symbol of arity \(l\).  If \(\kappa^{-1}(\ast)\) is empty, call
\((R,\kappa)\) a \emph{ground pattern}.  Fix a map \(\mathit{val}\) from
ground patterns to \(\{\top,\bot\}\).  For a formula \(\theta\) built from the
symbols of \(\varphi^{\flat}\) and a map \(\sigma\) from a finite set of
variables containing \(\mathrm{free}(\theta)\) to \(\{\ast,1,\ldots,p\}\), put
\([c_j]_\sigma\defeq j\) and \([x]_\sigma\defeq\sigma(x)\), and define
\(\mathrm{sp}_\sigma(\theta)\) by recursion.
\begin{itemize}
\item \(\mathrm{sp}_\sigma R(t_1,\ldots,t_r)\) is
  \(R_\kappa(t_{p_1},\ldots,t_{p_l})\) for \(\kappa(i)\defeq[t_i]_\sigma\) when
  \(\kappa^{-1}(\ast)\) is nonempty, and \(\mathit{val}(R,\kappa)\) otherwise.
\item \(\mathrm{sp}_\sigma(t=t')\) is \(t=t'\) if \([t]_\sigma=[t']_\sigma=\ast\);
  \(\top\) if \([t]_\sigma=[t']_\sigma\neq\ast\); and \(\bot\) otherwise.
\item \(\mathrm{sp}_\sigma\) commutes with \(\neg,\land,\lor,\to,\leftrightarrow\).
\item For \(\theta=(\exists\bar y\,.\,\alpha)\chi\), let \(\sigma'\) range over
  the maps that are defined on \(\mathrm{dom}(\sigma)\cup\bar y\) and agree
  with \(\sigma\) outside \(\bar y\), and let \(\bar y_{\ast}\) be the subtuple of
  the \(y\in\bar y\) with \(\sigma'(y)=\ast\).  Then
  \(\mathrm{sp}_\sigma(\theta)\) is the disjunction over all \(\sigma'\) of
  \((\exists\bar y_{\ast}\,.\,\mathrm{sp}_{\sigma'}\alpha)\mathrm{sp}_{\sigma'}\chi\)
  if \(\bar y_{\ast}\) is nonempty, and of
  \(\mathrm{sp}_{\sigma'}\alpha\land\mathrm{sp}_{\sigma'}\chi\) otherwise.
\item For \(\theta=(\forall\bar y\,.\,\alpha)\chi\) it is the conjunction over
  the same \(\sigma'\) of
  \((\forall\bar y_{\ast}\,.\,\mathrm{sp}_{\sigma'}\alpha)\mathrm{sp}_{\sigma'}\chi\)
  if \(\bar y_{\ast}\) is nonempty, and of
  \(\mathrm{sp}_{\sigma'}\alpha\to\mathrm{sp}_{\sigma'}\chi\) otherwise.
\end{itemize}
By induction, \(\mathrm{free}(\mathrm{sp}_\sigma\theta)\) consists of variables
\(x\in\mathrm{free}(\theta)\) with \(\sigma(x)=\ast\): an atom keeps exactly its
variables with value \(\ast\), and a variable of \(\bar y\) outside
\(\bar y_{\ast}\) is replaced by a constant index and does not occur.  In the
quantifier clauses with \(\bar y_{\ast}\) nonempty, \(\mathrm{sp}_{\sigma'}\alpha\)
is a relation atom \(R_\kappa(\bar w)\), because the variables of
\(\bar y_{\ast}\) occur in \(\alpha\) with value \(\ast\); its variables are the
variables of \(\alpha\) with value \(\ast\), which include \(\bar y_{\ast}\)
and every free variable of \(\mathrm{sp}_{\sigma'}\chi\), since
\(\mathrm{free}(\chi)\subseteq\mathrm{free}(\alpha)\).  So, reading \(\top\)
and \(\bot\) as additional atoms that never serve as guards, rule~(3) of
\cref{def:p02-w02-p02-035-gf} applies and \(\mathrm{sp}_\sigma\theta\) is
guarded and contains no constants.

The truth values are then removed bottom-up by replacements that preserve
equivalence and free variables: the Boolean rules such as
\(\top\land\chi\rightsquigarrow\chi\), \(\bot\land\chi\rightsquigarrow\bot\) and
\(\neg\top\rightsquigarrow\bot\), together with
\((\exists\bar y\,.\,\alpha)\bot\rightsquigarrow\bot\),
\((\exists\bar y\,.\,\alpha)\top\rightsquigarrow(\exists\bar y\,.\,\alpha)\alpha\),
\((\forall\bar y\,.\,\alpha)\top\rightsquigarrow\top\) and
\((\forall\bar y\,.\,\alpha)\bot\rightsquigarrow(\forall\bar y\,.\,\alpha)\neg\alpha\).
The result for \(\mathrm{sp}_\emptyset(\varphi^{\flat})\) is \(\top\), \(\bot\), or a
guarded sentence without truth values.  Let \(\varphi^{\mathit{val}}\) be that
sentence, or \((\exists x\,.\,Ux)Ux\) for a fresh unary \(U\) in the case
\(\top\), or \((\exists x\,.\,Ux)\neg Ux\) in the case \(\bot\).

\emph{Step 5: correctness of the split.}  Let \(\mathfrak A\) be a structure
for \(\varphi^{\flat}\) in which \(c_1,\ldots,c_p\) have pairwise distinct
values and \(A_0\defeq A\setminus\{c_1^{\mathfrak A},\ldots,c_p^{\mathfrak A}\}\)
is nonempty, and let \(\mathfrak B\) be a structure with universe \(A_0\) for
all symbols \(R_\kappa\).  For \(\kappa\) as above and \(\bar u\in A_0^{l}\),
let \(\langle\kappa,\bar u\rangle\) be the \(r\)-tuple whose \(i\)-th entry is
\(c_{\kappa(i)}^{\mathfrak A}\) if \(\kappa(i)\neq\ast\) and \(u_h\) if
\(i=p_h\).  Say that \(\mathfrak A\), \(\mathfrak B\) and \(\mathit{val}\)
\emph{correspond} if \(\langle\kappa,\bar u\rangle\in R^{\mathfrak A}\) iff
\(\bar u\in R_\kappa^{\mathfrak B}\) for all \(R\), all \(\kappa\) with a place
\(\ast\) and all \(\bar u\), and
\((c_{\kappa(1)}^{\mathfrak A},\ldots,c_{\kappa(r)}^{\mathfrak A})\in R^{\mathfrak A}\)
iff \(\mathit{val}(R,\kappa)=\top\) for every ground pattern.

\emph{Truth lemma.}  Suppose they correspond.  For \(\theta\) and \(\sigma\) as
in Step 4 and an assignment \(s\) of the variables \(x\) with
\(\sigma(x)=\ast\) into \(A_0\), let \(\hat s(x)\defeq s(x)\) if
\(\sigma(x)=\ast\) and \(\hat s(x)\defeq c_{\sigma(x)}^{\mathfrak A}\)
otherwise.  Then \(\mathfrak A\entails\theta[\hat s]\) iff
\(\mathfrak B\entails\mathrm{sp}_\sigma\theta[s]\), with \(\top\) and \(\bot\)
read as true and false.  For an atom \(R(t_1,\ldots,t_r)\), the value of
\(t_i\) under \(\hat s\) is \(c_{\kappa(i)}^{\mathfrak A}\) if
\(\kappa(i)\neq\ast\) and \(s(t_i)\) if \(\kappa(i)=\ast\), so the value tuple
is \(\langle\kappa,(s(t_{p_1}),\ldots,s(t_{p_l}))\rangle\), or the tuple of
constants for a ground pattern, and correspondence gives the claim.  For an
equality, two terms with value \(\ast\) take values \(s(t),s(t')\); two terms
with indices \(j,j'\) take values \(c_j^{\mathfrak A},c_{j'}^{\mathfrak A}\),
equal iff \(j=j'\) by distinctness; and a term with value \(\ast\) takes a
value in \(A_0\), which differs from every \(c_j^{\mathfrak A}\).  Boolean
steps are immediate.  For \(\theta=(\exists\bar y\,.\,\alpha)\chi\),
\(\mathfrak A\entails\theta[\hat s]\) iff some \(\bar b\) over \(A\) gives
\(\mathfrak A\entails(\alpha\land\chi)[\hat s,\bar y\mapsto\bar b]\).  A tuple
\(\bar b\) determines \(\sigma'\), with \(\sigma'(y)\defeq\ast\) if
\(b_y\in A_0\) and \(\sigma'(y)\defeq j\) for the unique \(j\) with
\(b_y=c_j^{\mathfrak A}\) otherwise, and the assignment \(s'\) that agrees with
\(s\) on the variables outside \(\bar y\) with value \(\ast\) and sends each
\(y\in\bar y_{\ast}\) to \(b_y\); then \(\hat s'\) is \(\hat s\) with
\(\bar y\mapsto\bar b\).  Every pair \((\sigma',s')\) of this form arises from
exactly one \(\bar b\).  By the induction hypothesis,
\(\mathfrak A\entails\theta[\hat s]\) iff for some \(\sigma'\) and some values
of \(\bar y_{\ast}\) in \(A_0\),
\(\mathfrak B\entails(\mathrm{sp}_{\sigma'}\alpha\land\mathrm{sp}_{\sigma'}\chi)[s']\),
which is the truth of the disjunct for \(\sigma'\), as \(A_0\) is the universe
of \(\mathfrak B\).  For \(\theta=(\forall\bar y\,.\,\alpha)\chi\) the same
correspondence between tuples \(\bar b\) and pairs \((\sigma',s')\) shows that
\(\alpha\to\chi\) holds for all \(\bar b\) iff
\(\mathrm{sp}_{\sigma'}\alpha\to\mathrm{sp}_{\sigma'}\chi\) holds for all
\(\sigma'\) and all values of \(\bar y_{\ast}\) in \(A_0\), which is the truth
of the conjunction.

\emph{Claim.}  \(\varphi_\equiv\) has a model in which its constants have
pairwise distinct values and some element is the value of no constant iff
\(\varphi^{\mathit{val}}\) is satisfiable for some \(\mathit{val}\).  Suppose
\(\varphi_\equiv\) has such a model; by Step 3, so does \(\varphi^{\flat}\),
call it \(\mathfrak A\).  Let \(\mathfrak B\) have universe \(A_0\) and
\(R_\kappa^{\mathfrak B}\defeq\{\bar u:\langle\kappa,\bar u\rangle\in R^{\mathfrak A}\}\),
and let \(\mathit{val}\) record which ground tuples lie in the relations of
\(\mathfrak A\).  They correspond, so
\(\mathfrak B\entails\mathrm{sp}_\emptyset(\varphi^{\flat})\); the
simplification is therefore not \(\bot\), and \(\mathfrak B\), with \(U\)
interpreted as \(A_0\) in the case \(\top\), satisfies
\(\varphi^{\mathit{val}}\).  Conversely let
\(\mathfrak B\entails\varphi^{\mathit{val}}\).  The case \(\bot\) is
impossible, since \((\exists x\,.\,Ux)\neg Ux\) has no model, so
\(\mathfrak B\entails\mathrm{sp}_\emptyset(\varphi^{\flat})\).  Interpret
every symbol \(R_\kappa\) absent from \(\varphi^{\mathit{val}}\) as empty.  Let
\(A\defeq B\cup\{d_1,\ldots,d_p\}\) with new distinct \(d_j\), and
\(c_j^{\mathfrak A}\defeq d_j\).  For \(\bar a\in A^{r}\) let
\(\kappa_{\bar a}(i)\defeq j\) if \(a_i=d_j\) and \(\ast\) otherwise, and put
\(\bar a\in R^{\mathfrak A}\) iff either \(\kappa_{\bar a}\) has places
\(p_1<\cdots<p_l\) with value \(\ast\) and
\((a_{p_1},\ldots,a_{p_l})\in R_{\kappa_{\bar a}}^{\mathfrak B}\), or it has none
and \(\mathit{val}(R,\kappa_{\bar a})=\top\).  For \(\bar u\) over \(B\),
\(\kappa_{\langle\kappa,\bar u\rangle}=\kappa\), so the three correspond, with
\(A_0=B\) nonempty.  By the truth lemma \(\mathfrak A\entails\varphi^{\flat}\),
and by Step 3 \(\varphi_\equiv\) has a model with distinct constant values on
the same universe, in which each \(b\in B\) is the value of no constant.

\emph{Step 6: the bounds.}  The sentence \(\varphi^{\mathit{val}}\) is
constant-free.  Its variables are among those of \(\varphi^{\flat}\) or are
the single variable \(x\), and each \(R_\kappa\) has arity at most that of
\(R\).  So if \(\varphi\) has at most \(k\) variables, \(\varphi^{\mathit{val}}\)
has at most \(q\); and if \(\varphi\) has arities at most \(k\),
\(\varphi^{\mathit{val}}\) has arities at most \(q\).  For a conjunct
\((\forall\bar x\,.\,\alpha)(\exists\bar y\,.\,\beta)\theta\), the tuples
\(\bar x\) and \(\bar y\) list distinct variables of the atom \(\beta\), so
\(\lvert\bar x\rvert+\lvert\bar y\rvert\leq q\); for the other two shapes the
bound tuple lists the variables of one atom.  Hence
\(\mathrm{sp}_\emptyset\) of a conjunct consists of at most \((p+1)^{q}\)
translated copies of its parts.  A translated part has the same number of
symbols as the part, or fewer, and each new symbol \(R_\kappa\) is written with
\(O(r\log(p+2))\) bits, so \(\lvert\varphi^{\mathit{val}}\rvert\) is at most
\((n+1)^{q}\) times a polynomial in \(n\), and the simplification does not
lengthen it.  It is computed in time polynomial in that length.  By
\cref{thm:p02-w02-p02-035-gf-upper} for constant-free input with at most
\(q\) variables or with arities at most \(q\), satisfiability of
\(\varphi^{\mathit{val}}\) is decided in time \(2^{n^{O(1)}}\).

The procedure tries every equivalence relation \(\equiv\), of which there are
at most \(p^{p}\leq2^{n\log n}\), and every restriction of \(\mathit{val}\) to
the ground patterns that occur in \(\mathrm{sp}_\emptyset(\varphi^{\flat})\).
An atom of \(\varphi^{\flat}\) has at most \(q\) distinct variables, so it
yields at most \((p+1)^{q}\) patterns, and there are at most
\(g\defeq\lvert\varphi^{\flat}\rvert(p+1)^{q}\) ground patterns, a polynomial
in \(n\); only these affect \(\varphi^{\mathit{val}}\).  It accepts iff some
\(\varphi^{\mathit{val}}\) is satisfiable.  By Steps 1 and 2 and the claim,
this happens iff \(\varphi\) is satisfiable.  The total time is at most
\(2^{n\log n}\cdot2^{g}\cdot2^{n^{O(1)}}=2^{n^{O(1)}}\).
\end{proof}
\end{annexitem}

\begin{annexitem}{annex:p02-w02-p02-035-gf-hard}{thm:p02-w02-p02-035-gf-hard}{Guarded satisfiability is 2EXPTIME-complete}
\begin{proof}
\emph{Acceptance.}  A configuration \(C=(q,r,y)\) of \(M\) consists of a
state, a head cell \(r<2^{s}\) and an inscription \(y\in\Gamma^{2^{s}}\); for
non-final \(q\), its \(i\)-th successor \(C_i\) applies
\(\delta_i(q,y_r)=(q',\sigma',\mu)\), giving \((q',r+\mu,y')\) with
\(y'_r=\sigma'\) and \(y'\) otherwise equal to \(y\).  Player \(\exists\)
wins from \(C\) if \(C\) is accepting, or existential with a winning
successor, or universal with both successors winning, and \(M\) accepts
\(w\) iff \(\exists\) wins from the initial configuration
\(C_{\mathrm{in}}=(q_0,0,w\square^{2^{s}-n})\) \cite[sect.~4.1]{Gradel1999}.
Since every computation halts and each configuration has at most two
successors, the configurations reachable from a reachable \(C\) form a
finite tree by K\"onig's lemma; let \(h(C)\) be its height.  Successors have
smaller height, so the recursion above is well founded.  By the machine
model, every reachable configuration with \(\mu=1\) has \(r+1<2^{s}\), and
with \(\mu=-1\) has \(r\geq1\).  We take \(p(n)\geq n\), so the input fits.

\emph{Form and size.}  Every quantifier of \(\mathrm{ACCEPT}_{M,w}\) carries
a relation-atom guard that contains all variables free in its scope: \(Rab\),
\(Tab\), \(U\bar xab\), \(S_iabcd\) and \(E\bar uab\bar vcd\) list the
relevant variables.  No equality and no constant occurs.  With \(s=p(n)\)
and \(M\) fixed, the sentence has \(O(s^2+ns)\) symbols and is produced in
polynomial time.

\emph{Facts at nodes.}  Let \(\mathfrak A\entails\mathrm{ACCEPT}_{M,w}\) and
\(\mathfrak A\entails Tab\).  Then \(Na\) and \(\neg Nb\), so \(a\neq b\), and
each \(r<2^{s}\) has its code \(\hat r(a,b)\).  For
\(\bar x,\bar x'\in\{a,b\}^{s}\), \(\mathit{succ}(\bar x,\bar x')\) holds iff
the number coded by \(\bar x'\) is one more than that coded by \(\bar x\):
the disjunct for \(i\) says that bits below \(i-1\) go from \(1\) to \(0\), bit
\(i-1\) goes from \(0\) to \(1\), and higher bits agree.  Likewise
\(\mathit{diff}(\bar x,\bar x')\) holds iff the coded numbers differ.  The
second and third conjuncts give \(U\bar xab\) for every
\(\bar x\in\{a,b\}^{s}\), by induction on the number of entries equal to
\(b\).  If also \(S_iabcd\) and \(Tcd\), call \(\bar v\in\{c,d\}^{2s}\)
\emph{matched} to \(\bar u\in\{a,b\}^{2s}\) if \(v_j=c\) iff \(u_j=a\); the
fifth conjunct gives \(E\bar uab\bar vcd\) for all matched pairs, again by
induction on the number of entries equal to \(b\).

Say that \(\mathfrak A\) \emph{imposes} \(C=(q,r,y)\) at \((a,b)\) if
\(Tab\), \(Q_qab\), \(Z\hat r(a,b)ab\), and \(P_{y_j}\hat j(a,b)ab\) for every
\(j<2^{s}\).

\emph{Step lemma.}  Let \(C\) be reachable and non-final, imposed at
\((a,b)\), and let \(S_iabcd\), \(Tcd\) and
\(\mathfrak A\entails\mathrm{NEXT}_i(a,b,c,d)\).  Then \(C_i\) is imposed at
\((c,d)\).  Write \(C=(q,r,y)\) and \(\delta_i(q,y_r)=(q',\sigma',\mu)\).  For
\(j\neq r\), instantiate the first conjunct of \(\mathrm{NEXT}_i\) at
\(\bar x=\hat j(a,b)\), \(\bar x'=\hat r(a,b)\), \(\bar y=\hat j(c,d)\),
\(\bar y'=\hat r(c,d)\); the guard holds because the pair is matched,
\(Z\bar x'ab\) and \(\mathit{diff}(\bar x,\bar x')\) hold, so
\(P_{y_j}\hat j(c,d)cd\).  For the head cell, instantiate the conjunct for
\((q,y_r)\) at \(\bar x=\hat r(a,b)\), \(\bar x'=\widehat{r+\mu}(a,b)\) and the
matched \(\bar y,\bar y'\); here \(r+\mu\) is a cell because \(C\) is
reachable, and \(\gamma_\mu\) holds.  This gives \(Q_{q'}cd\),
\(P_{\sigma'}\hat r(c,d)cd\) and \(Z\widehat{r+\mu}(c,d)cd\).  So \(C_i\) is
imposed at \((c,d)\).

\emph{Satisfiable implies acceptance.}  We show by induction on \(h(C)\)
that \(\exists\) wins from every reachable \(C\) imposed at some node.  If
\(C\) is final, it is not rejecting, by the last conjunct, so it is
accepting.  If \(C\) is existential, \(\mathrm{EXIST}(a,b)\) holds; the fourth
conjunct gives \(c,d\) with \(S_1abcd\) and \(Tcd\), and the seventh gives
\(\mathrm{NEXT}_1\) or \(\mathrm{NEXT}_2\) at \((a,b,c,d)\).  So by the step
lemma some \(C_i\) is imposed at \((c,d)\), and \(\exists\) wins from it by
induction.  If \(C\) is universal, the fourth conjunct gives
\((c_i,d_i)\) with \(S_iabc_id_i\) and \(Tc_id_i\) for \(i=1,2\), and the
eighth gives \(\mathrm{NEXT}_i\); so both successors are imposed and win.
Finally, the first conjunct gives \((a,b)\) with \(Rab\land Tab\), and the
sixth imposes \(C_{\mathrm{in}}\) there: cells \(j\geq n\) are covered by
\(U\hat j(a,b)ab\) and \(\mathit{diff}(\hat r,\hat j)\) for \(r<n\).  So
\(M\) accepts \(w\).

\emph{Acceptance implies satisfiable.}  Let \(\mathcal S\) be the tree of a
winning strategy of \(\exists\) from \(C_{\mathrm{in}}\): each node \(\nu\)
carries a configuration \(C(\nu)\), an existential node has one child
carrying a chosen successor \(C_{k(\nu)}\), a universal node has two
children carrying \(C_1\) and \(C_2\), and every leaf is accepting.  It is
finite.  Take distinct elements \(a_\nu,b_\nu\) for the nodes and interpret:
\(T=\{(a_\nu,b_\nu)\}\); \(R\) the root pair; \(N=\{a_\nu\}\);
\(U=\{(\bar x,a_\nu,b_\nu):\bar x\in\{a_\nu,b_\nu\}^{s}\}\);
\(S_i\) contains \((a_\nu,b_\nu,a_\mu,b_\mu)\) when \(\mu\) is the child of a
universal \(\nu\) carrying \(C_i\), or the only child of an existential
\(\nu\); \(E\) contains \((\bar u,a_\nu,b_\nu,\bar v,a_\mu,b_\mu)\) for such
quadruples and matched \(\bar u,\bar v\); \(Q_q\) contains
\((a_\nu,b_\nu)\) iff \(C(\nu)\) has state \(q\); \(Z\) contains
\((\hat r(a_\nu,b_\nu),a_\nu,b_\nu)\) for the head cell \(r\) of \(C(\nu)\);
and \(P_\sigma\) contains \((\hat j(a_\nu,b_\nu),a_\nu,b_\nu)\) iff cell \(j\)
of \(C(\nu)\) holds \(\sigma\).  Every node pair imposes exactly its own
configuration, and only through code tuples.

The first three conjuncts hold by construction.  The fourth holds because a
node whose state is existential or universal is not a leaf.  The fifth holds
because \(E\) contains exactly the matched tuples for each successor
quadruple, and the set of matched pairs contains \((\bar a,\bar c)\) and is
closed under the substitutions.  The sixth holds at the root because
\(C_{\mathrm{in}}\) is imposed there and the \(U\)-tuples are exactly the
codes.  For the seventh and eighth, let \(\mu\) be an \(S_i\)-child of
\(\nu\) with \(C(\mu)\) the \(k\)-th successor of \(C(\nu)\); then
\(\mathrm{NEXT}_k\) holds at the quadruple.  The \(E\)-tuples there are
exactly matched codes; if \(Z\bar x'ab\) and \(\mathit{diff}(\bar x,\bar x')\),
then \(\bar x\) codes a cell other than the head, whose symbol is unchanged
in \(C(\mu)\).  In a transition conjunct the antecedent forces \(q\), the
head cell and its symbol to be those of \(C(\nu)\), and \(\gamma_\mu\)
forces \(\bar x'\), and so \(\bar y'\), to code the new head cell; the
consequent then describes \(C(\mu)\).  For an existential \(\nu\) the
seventh conjunct is met by \(\mathrm{NEXT}_{k(\nu)}\), and the eighth is
vacuous; for a universal \(\nu\) the \(S_i\)-child carries \(C_i\), so
\(\mathrm{NEXT}_i\) holds, and the seventh is vacuous.  The ninth holds since
no node of \(\mathcal S\) is rejecting.  So
\(\mathrm{ACCEPT}_{M,w}\) has a model.

\emph{\(2\mathsf{EXPTIME}\)-completeness.}  Let
\(L\in\mathsf{DTIME}(2^{2^{n^k}})\).  By premise (i) with
\(s(n)=2^{n^k}\geq\log n\), \(L\) is decided by an alternating machine in
space \(2^{n^k}\), so by premise (ii) by a machine \(M\) as above with
\(p(n)\defeq n^k+n\).  The map \(w\mapsto\mathrm{ACCEPT}_{M,w}\) is a
polynomial-time reduction of \(L\) to satisfiability of equality-free
guarded sentences.  Membership in \(2\mathsf{EXPTIME}\) for the class with
equality and constants is \cref{thm:p02-w02-p02-035-gf-upper}.

\emph{\(\mathsf{EXPTIME}\)-hardness in \(\mathrm{GF}_2\).}  Let
\(L\in\mathsf{DTIME}(2^{n^k})\).  By premises (i) and (ii), \(L\) is decided
by a machine \(M\) as above within space \(S\defeq n^k+n\).  Use unary \(T\),
\(R\), \(Q_q\), \(Z_r\) and \(P_{\sigma,r}\) for \(r<S\), and binary
\(S_1,S_2\).  Let \(\mathrm{EXIST}(x)\), \(\mathrm{UNIV}(x)\) and
\(\mathrm{REJECT}(x)\) be the disjunctions of \(Q_qx\) over the respective
states, and let \(\mathrm{NX}_i(x,y)\) be the conjunction of
\((Z_{r'}x\land P_{\sigma,r}x)\to P_{\sigma,r}y\) for all \(\sigma\) and all
\(r\neq r'\) below \(S\), and of
\((Q_qx\land Z_rx\land P_{\sigma,r}x)\to(Q_{q'}y\land P_{\sigma',r}y\land Z_{r+\mu}y)\)
for all non-final \(q\), all \(\sigma\), and all \(r\) with \(0\leq r+\mu<S\),
where \(\delta_i(q,\sigma)=(q',\sigma',\mu)\).  The sentence is the conjunction
of
\[
\begin{gathered}
 (\exists x\,.\,Rx)Tx,\qquad
 (\forall x\,.\,Tx)\neg\mathrm{REJECT}(x),\\
 (\forall x\,.\,Tx)\bigl((\mathrm{EXIST}(x)\lor\mathrm{UNIV}(x))\to
  (\exists y\,.\,S_1xy)Ty\land(\exists y\,.\,S_2xy)Ty\bigr),\\
 (\forall x\,.\,Rx)\bigl(Q_{q_0}x\land Z_0x\land
  \textstyle\bigwedge_{r<n}P_{w_r,r}x\land\bigwedge_{n\leq r<S}P_{\square,r}x\bigr),\\
 (\forall xy\,.\,S_1xy)\bigl(\mathrm{EXIST}(x)\to\mathrm{NX}_1\lor\mathrm{NX}_2\bigr),\qquad
 \textstyle\bigwedge_{i=1,2}(\forall xy\,.\,S_ixy)\bigl(\mathrm{UNIV}(x)\to\mathrm{NX}_i\bigr).
\end{gathered}
\]
It uses only \(x\) and \(y\), relation symbols of arity at most \(2\), no
equality and no constants.  For fixed \(M\) it has \(O(S^2)\) conjuncts of
bounded size, each symbol name written with \(O(\log S)\) bits, so it is
computed in polynomial time.

Configurations, successors, the winning condition and the height \(h(C)\)
are as in the acceptance paragraph, with \(S\) cells in place of \(2^{s}\).
Say that an element \(e\) \emph{imposes} \(C=(q,r,y)\), \(y\in\Gamma^{S}\), if
\(Te\), \(Q_qe\), \(Z_re\) and \(P_{y_j,j}e\) for all \(j<S\).

\emph{Step lemma.}  Let \(C=(q,r,y)\) be reachable and non-final, imposed at
\(e\), and let \(S_iee'\), \(Te'\) and \(\mathfrak A\entails\mathrm{NX}_i(e,e')\).
Write \(\delta_i(q,y_r)=(q',\sigma',\mu)\).  For \(j\neq r\), the conjunct
\((Z_rx\land P_{y_j,j}x)\to P_{y_j,j}y\) of \(\mathrm{NX}_i\) applies, because
\(Z_re\) and \(P_{y_j,j}e\) hold, and gives \(P_{y_j,j}e'\).  Since \(C\) is
reachable, \(0\leq r+\mu<S\), so the transition conjunct for \(q\),
\(\sigma=y_r\) and \(r\) is present; its antecedent holds at \(e\), and it
gives \(Q_{q'}e'\), \(P_{\sigma',r}e'\) and \(Z_{r+\mu}e'\).  With \(Te'\), this
says that \(e'\) imposes \(C_i\).

\emph{Satisfiable implies acceptance.}  Let \(\mathfrak A\) be a model.  We show
by induction on \(h(C)\) that \(\exists\) wins from every reachable \(C\)
imposed at some element \(e\).  If \(C\) is final, it is not rejecting, by
the conjunct \((\forall x\,.\,Tx)\neg\mathrm{REJECT}(x)\), so it is accepting.
If \(C\) is existential, \(\mathrm{EXIST}(e)\) holds; the third displayed
conjunct gives \(e'\) with \(S_1ee'\) and \(Te'\), and the conjunct for
\(S_1\) and \(\mathrm{EXIST}\) gives \(\mathrm{NX}_1(e,e')\) or
\(\mathrm{NX}_2(e,e')\).  By the step lemma some successor \(C_i\) is imposed
at \(e'\), it has smaller height, and \(\exists\) wins from it by induction.
If \(C\) is universal, the third conjunct gives \(e_1,e_2\) with \(S_iee_i\)
and \(Te_i\), and the conjuncts for \(\mathrm{UNIV}\) give
\(\mathrm{NX}_i(e,e_i)\) for \(i=1,2\); so both successors are imposed and
\(\exists\) wins from both.  Finally, the first conjunct gives \(e\) with
\(Re\) and \(Te\), and the fourth gives \(Q_{q_0}e\), \(Z_0e\),
\(P_{w_r,r}e\) for \(r<n\) and \(P_{\square,r}e\) for \(n\leq r<S\), so \(e\)
imposes \(C_{\mathrm{in}}=(q_0,0,w\square^{S-n})\).  Hence \(M\) accepts
\(w\).

\emph{Acceptance implies satisfiable.}  Let \(\mathcal S\) be the finite tree of
a winning strategy of \(\exists\) from \(C_{\mathrm{in}}\), as in the
corresponding paragraph above, and take one element \(e_\nu\) for each node
\(\nu\).  Interpret \(T\) as the set of all \(e_\nu\); \(R\) as \(\{e_\nu\}\)
for the root \(\nu\); \(Q_q\) as the \(e_\nu\) whose \(C(\nu)\) has state
\(q\); \(Z_r\) as the \(e_\nu\) whose \(C(\nu)\) has head cell \(r\);
\(P_{\sigma,r}\) as the \(e_\nu\) whose \(C(\nu)\) has \(\sigma\) in cell
\(r\); and \(S_i\) as the pairs \((e_\nu,e_{\nu'})\) such that \(\nu'\) is the child
of a universal \(\nu\) carrying \(C_i\), or the only child of an existential
\(\nu\).  Each \(e_\nu\) then satisfies exactly one \(Q_q\), exactly one
\(Z_r\), and for each \(r\) exactly one \(P_{\sigma,r}\), namely those of
\(C(\nu)\).  The first and fourth conjuncts hold at the root, which carries
\(C_{\mathrm{in}}\).  The second holds because no node of \(\mathcal S\) is
rejecting.  The third holds because a node whose state is existential or
universal is not a leaf: an existential node has a child, which is an
\(S_1\)- and an \(S_2\)-successor, and a universal node has one child for each
\(i\).  For the last two conjuncts, let \(S_ie_\nu e_{\nu'}\), where
\(C(\nu')\) is the \(i'\)-th successor of \(C(\nu)=(q,r,y)\), and let
\(\delta_{i'}(q,y_r)=(q',\sigma',\mu)\).  Then \(\mathrm{NX}_{i'}(e_\nu,e_{\nu'})\)
holds.  In a conjunct \((Z_{r''}x\land P_{\sigma,j}x)\to P_{\sigma,j}y\) with
\(j\neq r''\), the antecedent forces \(r''=r\) and \(\sigma=y_j\); cell
\(j\neq r\) is unchanged in \(C(\nu')\), so the consequent holds.  In a
transition conjunct the antecedent forces the state \(q\), the head cell
\(r\) and the symbol \(y_r\) of \(C(\nu)\); the consequent then states the
new state \(q'\), the new symbol \(\sigma'\) in cell \(r\) and the new head
cell \(r+\mu\), which are those of \(C(\nu')\).  For an existential \(\nu\)
the conjunct for \(S_1\) is met by \(\mathrm{NX}_{i'}\) with \(i'=k(\nu)\), and
the conjuncts for \(\mathrm{UNIV}\) are vacuous; for a universal \(\nu\) the
\(S_i\)-child carries \(C_i\), so \(i'=i\) and \(\mathrm{NX}_i\) holds, and
the conjunct for \(\mathrm{EXIST}\) is vacuous.  So the sentence has a
model, and it is satisfiable iff \(M\) accepts \(w\).
This sentence lies in \(\mathrm{GF}_{k'}\) and has arities at most \(k'\) for
every \(k'\geq2\), which gives the stated hardness.  Together with
\cref{thm:p02-w02-p02-035-gf-upper} for constant-free input and
\cref{thm:p02-w02-p02-035-gf-bounded-constants} for input with constants,
bounded-variable and bounded-arity guarded satisfiability are
\(\mathsf{EXPTIME}\)-complete; the hard sentences contain no constants, so
hardness holds for both input languages.
\end{proof}
\end{annexitem}
